%%%%%%%%%%%%%%%%%%%%%%%%%%%%%%%%%%%%%%%%%%%
\section{Tame fundamental groups of rigid spaces}
\label{s:fg-pi1-rig}
%%%%%%%%%%%%%%%%%%%%%%%%%%%%%%%%%%%%%%%%%%%


In this section, we define and study the tame fundamental group $\pitame(X/K)$ of a connected rigid-analytic space relative to a base non-archimedean field $K$ (see \S\ref{ss:adic-spaces} for our conventions and basic facts regarding rigid geometry). A finite \'etale cover $f\colon Y\to X$ is said to be tame relative to $K$ if for every $y\in Y$, the extension $k(y)/k(f(x))$ is tamely ramified with respect to every valuation subring $V$ of $k(y)$ containing $K^\circ$ (see \cref{def:tame-rig}). As explained in \S\ref{sec:root-triples}, if $X$ is separated and taut, such pairs $(y,V)$ are the points of the universal compactification $\overline{Y} = \Spa(Y, K)$, and $Y\to X$ is tame relative to $K$ if and only if for the induced finite \'etale map $\overline{Y}\to \overline{X}$ the extensions of residue fields are tamely ramified (see \cref{prop:tame-S-compactification}). Finite \'etale covers of $X$ which are tame relative to $K$ form a Galois category with fundamental group~$\pitame(X/K)$.

Our main technical result (\cref{cor:Abhy-infty-equiv}, proved in \S\ref{ss:log-abhyankar}--\ref{ss:tameness-comparison}) states that if $K$ is either discretely valued or algebraically closed, then for a connected strictly semistable formal scheme $\fX$ over $K^\circ$ with generic fibre $X$ and special fibre $\fX_k$, we have an isomorphism
\[
    \pitame(X/K) \simeq \pitame(\fX^{\rm log}_k/k)
\]
where $\fX^{\rm log}$ is the log formal scheme obtained by endowing $\fX$ with the log structure induced by the generic fibre (called the ``standard log structure'' below), and $\fX^{\rm log}_k$ is the special fibre endowed with the log structure induced from $\fX^{\rm log}$. Then $\fX^{\rm log}_k$ is a log scheme of type $\typeVd$ over $k$ and the group on the right is the one studied in \S\ref{s:fg-pi1-log}. In particular, combined with \cref{thm:pi1tame log is fg} this implies immediately that $\pitame(X/K)$ is finitely generated if $\fX$ is quasi-compact and $k$ is algebraically closed. 

Using an auxiliary result about the behaviour of $\pitame(X/K)$ with respect to field extensions (\S\ref{ss:K-vs-C}), some descent assertions (\S\ref{ss:alterations-v-descent}), and uniformisation results of Gabber (\cref{prop:h-locally-regular}) and Temkin (\cref{thm:Temkin-unif}), we can drop the semistable reduction hypothesis, as well as the restrictions on $K$, and show that $\pitame(X/K)$ is finitely generated if $X$ is qcqs and if the tame Galois group $\pitame(K)$ is finitely generated (\cref{thm:fg-pi1-rig}). Employing \cref{cor:surgery-fp} we also show that $\pitame(X/K)$ is finitely presented under more restrictive assumptions (\cref{thm:fp-sometimes}). We conclude this section with a few examples (see \cref{ex:polydisc}--\cref{ex:helminck}).


%%%%%%%%%%%%%%%%%%%%%%%%%%%%%%%%%%%%%%%%%%
\subsection{The tame fundamental group of a rigid space}
\label{ss:tame-pi1-rig}

We give the definition of the tame fundamental group for arbitrary morphisms of adic spaces, before specialising to the case of rigid-analytic varieties. First, recall the following definition from \cite{Hubner2021:AdicTameSite}. 

\begin{defi}[{Absolute tameness}] \label{def:tame-rig-abs}
    Let $f\colon Y\to X$ be an \'etale morphism of adic spaces. We say that $f$ is {\bf tame} if for every $y\in Y$ with image $x\in X$,  the finite separable extension of valued fields 
    \[
        (k(y), k(y)^+)/(k(x), k(x)^+)
    \]
    is tamely ramified, see \S\ref{sec:tame extensions valued fields}.
\end{defi}

We extend this definition by making it relative to a base adic space $S$.

\begin{defi}[Relative tameness] \label{def:tame-rig}
    Let $X\to S$ be a morphism of adic spaces and let $Y\to X$ be a finite \'etale morphism. We say that $Y\to X$ is \textbf{tame relative to $S$} if for every root triple 
    \[
        y= (y^\circ, k(y)^+, y^\succ) \in \Spa(Y, S)
        \qquad \text{(see~\cref{defi:root-triple})}
    \]
    with image $x = (x^\circ, k(x)^+, x^\succ)\in \Spa(X, S)$, the finite separable extension of valued fields $(k(y), k(y)^+)/(k(x), k(x)^+)$ is tamely ramified. 
\end{defi}

\begin{rmks} \label{rmks:relative-tameness}
    \begin{enumerate}[(1)]
        \item Note that we defined tameness relative to~$S$ only for \emph{finite} étale morphisms. So with this notion we will be able to define a tame fundamental group of~$X$ relative to~$S$ (see \cref{def:tamepi1-rigid}) but we do not talk about a tame site of~$X$ relative to~$S$. We could define tameness relative to $S$ for a general étale morphism just as above but this is not a useful concept. For instance, $\Spa(Y,S) \to \Spa(X,S)$ fails to be étale if $Y \to X$ is étale but not finite. Moreover, many of the properties proved in \cref{lem:properties-tameness-relative-S} below are not satisfied if $Y\to X$ is not finite. 
        \item In the setting of \cref{def:tame-rig}, any point $y \in Y$ gives rise to the root triple $(y^\circ,k(y)^+,y^\succ)$ where~$y^\succ$ is the image of~$y$ in~$S$. So if $Y \to X$ is tame relative to~$S$, it is also (absolutely) tame. For tameness relative to~$S$ there might be additional tameness conditions as there can exist root triples that do not arise from points of~$Y$. 
        \item \label{enumitem:relative-tameness-test-square} A finite \'etale $Y\to X$ is tame relative to $S$ if and only if the following condition holds: for every point of $\Spa(X, S)$ represented by a square \labelcref{eqn:pp-test-square}, the pull-back of $Y$ to $\Spa(k, k^\circ)$ is a finite \'etale map. Since $\Spa(k, k^\circ)$ and $\Spa(k, k^+)$ have the same finite \'etale covers (namely, finite separable $k$-algebras), we obtain a finite \'etale cover of $\Spa(k, k^+)$, and we want it to be tame with respect to $k^+$. 
    \end{enumerate}
\end{rmks}

\begin{lem} \label{lem:properties-tameness-relative-S}
    Let $Y\to X\to S$ be maps of adic spaces, with $Y\to X$ finite \'etale. 
    \begin{enumerate}[(a)]
        \item \label{lemitem:sorite-tameness-relative-S1}
        If $X=S$, then $Y\to X$ is tame relative to $S$ if and only if it is tame (in the sense of \cref{def:tame-rig-abs}).
        \item \label{lemitem:sorite-tameness-relative-S2}
        Given a map $S\to T$, if $Y\to X$ is tame relative to $T$, then it is tame relative to $S$. The converse holds if $S\to T$ is vertically partially proper.
        \item 
        \label{lemitem:sorite-tameness-relative-S3}
        Given a surjective finite \'etale $Z\to Y$, the map $Z \to X$ is tame relative to~$S$ if and only if $Z\to Y$ and $Y\to X$ are tame relative to $S$. 
        \item 
        \label{lemitem:sorite-tameness-relative-S4}
        Given a commutative diagram
        \[ 
            \begin{tikzcd}
                Y'\ar[r] \ar[d] & Y \ar[d] \\
                X'\ar[r] \ar[d] & X \ar[d] \\
                S'\ar[r]  & S
            \end{tikzcd}
        \]
        where $Y' = Y\times_X X'$, if $Y\to X$ is tame relative to $S$, then $Y'\to X'$ is tame relative to~$S'$.
    \end{enumerate}    
\end{lem}

\begin{proof}
\labelcref{lemitem:sorite-tameness-relative-S1} In this case $\Spa(X, S) = X$ and $\Spa(Y,S) = Y$.
The first part of \labelcref{lemitem:sorite-tameness-relative-S2} is a special case of \labelcref{lemitem:sorite-tameness-relative-S4} (take  $X=X'$), and the second follows from the fact that $\Spa(X, S)\to \Spa(X, T)$ is a bijection if $S\to T$ is vertically partially proper (see \cref{prop:Spa-iso-vpp}).
Part \labelcref{lemitem:sorite-tameness-relative-S3} follows from the standard fact that for algebraic extensions of valued fields $M/L/K$, the extension $M/K$ is tame if and only if $M/L$ and $L/K$ are tame. Finally, part \labelcref{lemitem:sorite-tameness-relative-S4} is clear by the characterisation in \cref{rmks:relative-tameness}~\ref{enumitem:relative-tameness-test-square}.
\end{proof}

We can understand tameness relative to~$S$ in terms of the universal vertical compactification.

\begin{prop} \label{prop:tame-S-compactification}
    Let $Y\to X\to S$ be maps of adic spaces.
    Suppose that $Y\to X$ is finite \'etale and that $X\to S$ satisfies the assumptions of \cref{thm:Solodov} and \cref{thm:Solodov-comp}. Then the map $Y\to S$ also satisfies said assumptions, and the induced map (see \cref{prop:Spa-functoriality}) between universal vertical compactifications relative to $S$
        \[ 
            \bar{f}\colon \Spa(Y,S)\longrightarrow \Spa(X,S)
        \]
    is finite \'etale. The map $f$ is tame relative to $S$ if and only if the map $\bar{f}$ is tame (in the absolute sense of \cref{def:tame-rig-abs}).
\end{prop}

\begin{proof}
We first convince ourselves that~$\bar{f}$ is finite étale.
We start with the affinoid case.
If~$S$ and~$X$ are affinoid, so is~$Y$, and we can write
\[
 S = \Spa(R,R^+), \qquad X = \Spa(A,A^+), \qquad Y = \Spa(B,B^+),
\]
for complete Huber pairs $(R,R^+)$, $(A,A^+)$, and $(B,B^+)$.
Then
\[
 \Spa(X,S) = \Spa(A,R^+) = \Spa(A,A^+_{\min}), \qquad \Spa(Y,S) = \Spa(B,R^+) = \Spa(B,B^+_{\min}),
\]
(see \cref{prop:affinoid-compactification} and \cref{rmk:hypotheses-compactification}) where~$A^+_{\min}$ and~$B^+_{\min}$ are the smallest rings of integral elements in~$A$ and~$B$, respectively, containing~$R^+$.
Since $Y \to X$ is finite étale, $B$ is finite étale over~$A$ and~$B^+$ is the integral closure of~$A^+$ in~$B$.
The integral closure of~$A^+_{\min}$ in~$B$ is a ring of integral elements in~$B$ and must therefore equal~$B^+_{\min}$.
We conclude that~$\bar{f}$ is finite étale in the affinoid case.

In order to deduce the general case, we cover~$S$ by affinoid opens~$S_i$.
Then by \cref{cor:functoriality-compactification}, 
\[
 \Spa(X \times_S S_i,S_i) \simeq \Spa(X,S) \times_S S_i
\]
and these spaces form an open covering of $\Spa(X,S)$.
The restriction of~$\bar{f}$ to $\Spa(X \times_S S_i,S_i)$ is the natural map
\[
 \Spa(Y \times_S S_i,S_i) \la \Spa(X \times_S S_i,S_i)
\]
and since the property of being finite étale is local on the target, we are reduced to the case where~$S$ is affinoid.

Now we cover~$X$ by open affinoids~$X_i$.
We obtain a surjective family
\[
 \{\Spa(X_i,S) \la \Spa(X,S)\}_i
\]
of subspaces.
By \cref{lem:interiors-cover}, the open subspaces $\Spa(X_i,S)^\circ \subseteq \Spa(X_i,S)$ defined there comprise an open cover of $\Spa(X,S)$.
Using \cref{prop:comp-fibre-product} we identify the restriction of~$\bar{f}$ to $\Spa(X_i,S)$ with the morphism
\[
 \Spa(X_i \times_X Y,S) \longrightarrow \Spa(X_i,S).
\]
By the affinoid case this morphism is finite étale and stays so when restricted to~$\Spa(X_i,S)^\circ$.
As the open subsets~$\Spa(X_i,S)^\circ$ cover $\Spa(X,S)$, this shows that~$\bar{f}$ is finite étale.

The last assertion about~$f$ being tame relative to~$S$ if and only if~$\bar{f}$ is tame follows immediately by the description of the points of $\Spa(Y,S)$ as root triples.
\end{proof}

By a {\bf geometric point} of an adic space $X$ we shall mean a map of adic spaces $\Spa(L,L^+)\to X$ where $(L, L^+)$ is an affinoid field with $L$ separably closed, see \cite[\S 3.2]{AchingerLaraYoucis2023:GeometricArcsFundamental}. Note that the space $\Spa(L, L^+)$ might have more than one point.

\begin{lem} 
\label{lemma:adic-tame-galois}
    Let $X\to S$ be a morphism of adic spaces. Suppose that $X$ is connected and let $F_{\ov x}$ be the fibre functor associated to a geometric point  $\ov{x}$ of $X$. 

    Then $\FEtt_{X/S}$ is a Galois category with fibre functor $F_{\ov{x}}$. 
\end{lem}

\begin{proof}
This follows from the criterion of \cref{lem:Galois-subcategory-criterion} and \cref{lem:properties-tameness-relative-S}. 
\end{proof}

\begin{defi} \label{def:tamepi1-rigid}
\begin{enumerate}
    \item 
    In the situation of \cref{lemma:adic-tame-galois}, we denote by 
    $\pitame(X/S, \ov{x})$
    the fundamental group of the Galois category $\FEtt_{X/S}$ with base point $F_{\ov{x}}$, and call it the \textbf{tame fundamental group} of $X$ relative to $S$. 
    \item
    In the special case $X = S$ we call $\pitame(X, \ov{x}) = \pitame(X/X, \ov{x})$ the \textbf{(absolute) tame fundamental group} of $X$. The latter invariant has been introduced in \cite[\S9]{Hubner2021:AdicTameSite}.
    \item    
    If $X$ is a rigid-analytic space over a non-archimedean field $K$, we write $\pitame(X/K, \ov{x})$ for the tame fundamental group of $X$ relative to $\Spa(K)$. 
\end{enumerate}
\end{defi}

\begin{rmk} \label{rmk:Berkovich-covers}
    In \cite[\S6.3]{Berkovich:EtaleCohomology} Berkovich studies a variant of tame covers of Berkovich spaces over a non-archimedean field of residue characteristic $p > 0$.
    In the language of adic spaces, the (connected) covers he considers are (connected) finite étale covers $Y \to X$ whose Galois closure $Y' \to X \in \FEt_X$ satisfies the following: for any point $y' \in Y'$ with image $x \in X$, the degree $[\widehat{k(y')}:\widehat{k(x)}]$ is prime to~$p$.
    Every such cover is tame relative to~$K$ in our sense.
    However, a nontrivial strongly étale $p$-cover is tame relative to~$K$ but not tame in Berkovich's sense.
    Berkovich's covers form a Galois subcategory of $\FEtt_{X/K}$ by~\cref{lem:Galois-subcategory-criterion} and the corresponding fundamental group~$\pi_1^B(X/K)$ is a quotient of~$\pitame(X/K)$.
\end{rmk}

Since for $X\to S\to T$ we have full subcategories 
\[
    \FEtt_{X/T}\subseteq  \FEtt_{X/S}\subseteq \FEtt_X\subseteq \FEt_X,
\]
we obtain continuous surjections of profinite groups
\[ 
    \pi_1(X, \overline{x}) \surj \pitame(X, \overline{x}) \surj \pitame(X/S,\overline{x}) \surj \pitame(X/T,\overline{x}).
\]
In particular, for a rigid-analytic space $X$ over $K$, we obtain surjections
\begin{equation} \label{eqn:pi1comp-maps}
    \pi_1(X, \overline{x}) \surj \pitame(X,\overline{x}) \surj \pitame(X/K,\overline{x}).
\end{equation}
Neither of these maps is an isomorphism in general, as the following example shows. 

\begin{ex} \label{ex:unit-disc-Cp}
    Let $K = \mathbb{C}_p$ be a completed algebraic closure of $\mathbb{Q}_p$ and let 
    \[ 
        X = \Spa(K\langle T\rangle)
    \]
    be the affinoid unit disc over $K$. In the sequence of surjections \eqref{eqn:pi1comp-maps}, neither of the maps is an isomorphism, and only the right-most group is finitely generated (in fact trivial, see \cref{ex:polydisc}). 
    
    Let $\mathfrak{X} = \Spf(K^\circ \langle T\rangle)$ be the standard affine formal model of $X$ with special fibre $\mathbb{A}^1_k$ where $k=K^\succ$ is the residue field of $K$. Every finite \'etale map $Y_0\to \mathbb{A}^1_k$ lifts uniquely to a finite \'etale $\mathfrak{Y}\to \mathfrak{X}$. Its generic fibre $Y=\mathfrak{Y}_{\rm rig}\to \mathfrak{X}_{\rm rig}=X$ is a tame map \emph{relative to $X$}, and $Y$ is connected if and only if $Y_0$ is (since $\mathfrak{Y}$ is normal, cf.~\cite[Lemma A.2+Corollary A.16]{AchingerLaraYoucis2022}). We obtain a surjection
    \[ 
        \pitame(X,\overline{x}) \surj \pi_1(\mathbb{A}^1_k, \overline{x}).
    \]
    Since the target is not finitely generated, neither are $\pitame(X, \overline{x})$ and $\pi_1(X, \overline{x})$. As we shall prove later in \S\ref{ss:tameness-comparison}, the cover $Y\to X$ is tame \emph{relative to $K$} if and only if $Y_0 \to \mathbb{A}^1_k$ is tame relative to $k$ (in the sense of \cref{defi:tame cover of schemes X over S}), which happens only if $Y_0 \to\mathbb{A}^1_k$ is trivial.

    To be explicit, the finite \'etale covers
    \[ 
        Y_\lambda = \Spa(K\langle T, U\rangle/(U^p-U-\lambda T)) \longrightarrow \Spa(K\langle T\rangle) = X
    \]
    for $\lambda\in\mathbb{C}_p$ with $|\lambda|=1$ give rise to $\mathbb{F}_p$-torsors such that $Y_\lambda\neq Y_{\lambda'}$ if $|\lambda-\lambda'|=1$. They are all tame relative to $X$, so $\Hom(\pitame(X,\overline{x}),\bF_p)$ is infinite and therefore $\pitame(X,\overline{x})$ is not finitely generated. On the other hand, the $Y_\lambda \to X$ are not tame relative to $K$.  
    To see the latter, consider, as in \cref{ex:comp-disk}, the unique boundary point of the universal compactification 
    \[
        \Spa(X,K)\setminus X = \{\xi\}.
    \]
    Then the corresponding maximal point $\xi^\circ$ is the Gau\ss\ point. Its specialisation field $k(\xi^\circ)^\succ$ is the function field $k(T)$ of $\mathbb{A}^1_k$. If $v_\infty$ is the valuation on $k(T)$ corresponding to the point at infinity, then the valuation corresponding to $\xi$ is the composite of the two. In the cover $Y_\lambda \to X$ the Gau\ss\ point is inert with residue field extension $k(U)/k(T)$, the function field of the Artin-Schreier extension defined by $U^p-U - \bar \lambda T = 0$, where $0 \not=\bar \lambda \in k$ is the reduction of $\lambda$. It follows that $v_\infty$ is wildly ramified in this residual representation, hence the composite $\xi$ is also not tame by \cref{composition-tame}.

    In order to construct a finite \'etale $Y\to X$ which is not even tame relative to $X$, consider the finite \'etale cover of the disc described by the equation $U^p - p^{1/2}U  - T = 0$. Here the specialisation field extension at the Gau\ss\ point is purely inseparable.    
\end{ex}

%%%%%%%%%%%%%%%%%%%%%%%%%%%%%%%%%%%%%%%%%%%
\subsection{Behaviour with respect to field extensions} 
\label{ss:K-vs-C}

To deal with a non-algebraically closed base field $K$, we need to investigate how the tame fundamental group $\pitame(X/K)$ of a geometrically connected rigid space $X$ over $K$ compares to its geometric variant $\pitame(X_C/C)$ where $C$ is a completed algebraic closure of $K$. This relationship is slightly subtle, as the analogue of the fundamental sequence is not exact in general. Fortunately, it is always almost exact, as we show in \cref{prop:weak-fundamental-exact-seq} below. The case when $X$ is a point is covered by the following lemma. 

\begin{lem} \label{lem:structure of pitame na field}
    Let $(K,K^+)$ be a henselian valued field of residue characteristic $p \ge 0$.
    We fix an algebraic closure of $(K,K^+)$ giving rise to a geometric point $\bar x$ of $S=\Spa(K,K^+)$.
    We use the notation $\pitame(K)$ to denote the Galois group $\Gal(K^t/K)$.
    Then there is a canonical isomorphism
        \[
        \pitame(S,\bar x) \isomto \pitame(K).
        \]
\end{lem}

\begin{proof}
    The assertion follows immediately from the definition since the relevant space of root triples $\Spa(S,S) = S$ contains only coarsenings of the valuation defining $K^+$. 
\end{proof}

Let now $X$ be a geometrically connected qcqs rigid analytic space over a non-archimedean field $K$ and let $C$ be a completed algebraic closure of $K$. Pick a geometric point $\ov x$ of $X_C$. We are interested in the exactness of the following sequence of profinite groups 
\begin{equation} \label{eqn:fundamental-ex-seq}
    \begin{tikzcd}
        \pitame(X_C/C, \ov x)\ar[r] & \pitame(X/K,\ov x) \ar[r] & \pitame(K)\ar[r] & 1.
    \end{tikzcd}
\end{equation}
The subtlety surrounding this sequence which we hinted at in the first paragraph stems from the fact that for a wildly ramified finite separable extension $K'/K$, the induced finite \'etale morphism $X_{K'}\to X$ might be tame relative to $K$ (see \cref{ex:Raynaud inspired example,ex:moregeneral-twisted-nodal-curve} below). 
In order to address this issue, let us consider the full subcategory $\cC_{X/K}$ of $\FEt_K$ of all (spectra of) finite \'etale algebras $A/K$ such that 
\[
    X_A = X \times_{\Spa(K)} \Spa(A)\la X
\]
is tame relative to $K$. This contains $\FEtt_{K}$ but might a priori be larger if $X(K)=\emptyset$. The category $\cC_{X/K}$ is a Galois category whose fundamental group we denote by $\pi_1(\cC_{K/X})$. Clearly we have surjections
\[
\pi_1^\et(K) \surj \pi_1(\cC_{K/X}) \surj \pitame(K).
\]

\begin{prop} 
\label{prop:weak-fundamental-exact-seq}
    Let $X$ be a geometrically connected qcqs rigid analytic space over a non-archimedean field $K$. Let $C$ be the completion of an algebraic closure of $K$.  Let $\bar x$ be a geometric point of the base change $X_C$. We will also denote by $\bar x$ its image in $X$ and compute $\pi_1(\cC_{K/X}) \surj \pitame(K)$ with respect to the separable closure of $K$ in $\bar x$. 
    \begin{enumerate}[ref=(\arabic*)]
        \item 
        \label{propitem:weak-fes ct version}
        The sequence of profinite groups
        \begin{equation}
            \label{eq:homotopy exact sequence ct version}
            \begin{tikzcd}
                \pitame(X_C/C, \bar x)\ar[r] & \pitame(X/K, \bar x) \ar[r] & \pi_1(\cC_{K/X}) \ar[r] & 1           
            \end{tikzcd}
        \end{equation}        
        is exact.
        \item \label{propitem:finite kernel}
        The kernel of $\pi_1(\cC_{K/X}) \surj \pitame(K)$ is finite.
        \item 
        \label{propitem:weak-fes}
        If we assume furthermore that $X(K^t) \not= \varnothing$, then 
        $\pi_1(\cC_{K/X})\isomto \pitame(K)$ and
        \begin{equation}
            \label{eq:homotopy exact sequence}
            \begin{tikzcd}
                \pitame(X_C/C, \bar x)\ar[r] & \pitame(X/K, \bar x) \ar[r] & \pitame(K) \ar[r] & 1   
            \end{tikzcd}
        \end{equation}        
        is exact.
    \end{enumerate}
\end{prop}

\begin{proof}
    \labelcref{propitem:weak-fes ct version}
    The map $\pitame(X/K,\bar x) \surj \pi_1(\cC_{K/X})$ is surjective, because $X$ is assumed to be geometrically connected, hence $L \mapsto X_L$ maps connected objects to connected objects. 

    The composition $\pitame(X_C/C, \bar x) \to \pitame(X/K,\bar x) \surj \pi_1(\cC_{K/X})$ is the zero homomorphism, because for every $L$ in $\cC_{X/K}$ the cover $C \times_K X_L = \coprod_{L \inj C} X_C \to X_C$ is completely split. 

    We proceed as in \cite[X]{SGA1}. It remains to show that for every open subgroup $U \subseteq \pitame(X/K,\bar x)$ corresponding to a connected tame cover $X' \to X$ such that the projection  $X_C \to X$ admits a lift $X_C \to X'$, there is a separable extension $L/K$ such that $X' \simeq X_L$ as tame covers of $X$. This follows at once from
    \cite[Proposition~2.13]{deJong95:FundamentalGroupsofAnalytic}
    (which is stated for Berkovich spaces). Note that $L$ automatically belongs to $\cC_{X/K}$.
    
    \labelcref{propitem:finite kernel} The Galois category $\cC_{X/K}$ also makes sense if $X$ is not necessarily geometrically connected over $K$. Moreover, for a $K$-morphism $Y \to X$ of rigid analytic varieties, $\cC_{X/K}$ is contained in $\cC_{Y/K}$. In order to show \labelcref{propitem:finite kernel} we may replace $X$ by $Y$.
    
    Let $y = \Spa(K') \to X$ be a classical point with $K'/K$ finite. It thus suffices to show that there is a field extension $M$ of $K$ that is a finite extension of $K^t$ such that all fields $L$ from $\cC_{y/K}$ embed into $M$. A field $L$ belongs to $\cC_{y/K}$ if and only if $L \otimes_K K'$ is a product of tame extensions of $K'$. Being a subfield of those factors, $L$ is contained in $K'^t = K' K^t$, see  \cref{cor:pitame injective for fields}, which is finite over $K^t$. This verifies \labelcref{propitem:finite kernel} and assertion \labelcref{propitem:weak-fes} follows at once from the above proof if we pick $K' \subseteq K^t$.     
\end{proof}

\begin{cor} \label{cor:fg-alg-closed}
    In the situation of \cref{prop:weak-fundamental-exact-seq}, suppose that $\pitame(X_C/C,\bar x)$ and $\pitame(K)$ are finitely generated. Then $\pitame(X/K,\bar x)$ is finitely generated as well.
\end{cor}


\begin{proof}
By \cref{prop:weak-fundamental-exact-seq}~\labelcref{propitem:finite kernel} also 
$\pi_1(\cC_{K/X})$ is finitely generated. The claim thus follows from the short exact sequence from
\cref{prop:weak-fundamental-exact-seq}~\labelcref{propitem:weak-fes ct version}.
\end{proof}

\begin{rmk}
    The conclusions of \cref{prop:weak-fundamental-exact-seq} and \cref{cor:fg-alg-closed} also hold for a geometrically connected scheme of finite type over $K$ with tameness relative to $K^\circ$.
\end{rmk}

We finish this subsection by giving two examples of geometrically connected rigid spaces $X$ over $K$ for which $\pi_1(\cC_{X/K})\neq \pitame(K)$ and consequently for which \eqref{eqn:fundamental-ex-seq} is not exact. 

\begin{ex}[{due to Wittenberg based on ideas of Raynaud from  \cite[\S9]{Raynaud1970:SpecialisationFoncteurPicard}}] \label{ex:Raynaud inspired example}
    Let $K$ be a complete discretely valued field with residue field $k$ of characteristic $p>0$ and let $E$ be an elliptic curve over $K^\circ$ with a $p$-torsion point $x\in E(K^\circ)$ with nonzero image in $E(k)$. Translation by $x$ gives a free action of $\bZ/p\bZ$ on $E$ and an \'etale $\bZ/p\bZ$-torsor $E\to E/(\bZ/p\bZ)=F$.   
    Let $L/K$ be a wildly ramified $\mathbb{Z}/p\mathbb{Z}$-extension. We set
    \[
        X = \left(E\times_{\Spec(K^\circ)} \Spec(L^\circ)\right)\bigg/ (\bZ/p\bZ)
    \]
    with the quotient by the free diagonal action of $\bZ/p\bZ$, using the translation-by-$x$ action on $E$ and the Galois action on $L^\circ$. 
    So $X \to F$ is an $E/F$-twisted version of $F_{L^0} \to F$. Since the twisting $\bZ/p\bZ$-torsor $E\to F$ becomes trivial after pull-back to $E$, so does $X \to F$ and hence we have a cartesian square
    \[
        \begin{tikzcd}
            E_{L^0} \ar[r] \ar[d] & E \ar[d] \\
            X \ar[r] & F.
        \end{tikzcd}
    \]
    In particular, $E_{L^0}\to X$ is \'etale. On the other hand, the generic fibre of this map is the projection $E_L = X_L\to X_K$. To see this, note that $X_L=E_L\to F_K$ is a connected $(\bZ/p\bZ)^2$-torsor and that $X_K$ and $E_K$ are non-isomorphic degree $p$ intermediate coverings (as $X_K$ is a non-trivial twisted form of $E_K$), which implies that $E_L \simeq X_K\times_{F_K} E_K$, as desired.
    We conclude that the map of generic fibres $X_L\to X_K$ is tame relative to $K^\circ$, and that the induced map of rigid analytic spaces $X_L^{\rm an}=(X_K^{\rm an})_L\to X_K^{\rm an}$ is tame relative to $K$, showing that $\pi_1(\cC_{X_K^{\rm an}/K})\neq\pitame(K)$.
\end{ex}

\begin{ex}[Twisted nodal curve] \label{ex:moregeneral-twisted-nodal-curve}
Let $K$ be a non-archimedean field with a wildly ramified cyclic separable extension $L/K$. Let $\tau \in \Gal(L/K)$ be a generator. Let $X=\Spa(A)$ where
\[
    A= \{f\in K\langle T\rangle\,:\, f(0) = \tau(f(1))\}, 
\]
which is an affinoid curve with a single node with residue field $L$. We claim that $X_L\to X$ is tame relative to $K$. 
One computes that 
\[
    A\otimes_K L \simeq \left\{(f_\sigma) \in {\textstyle \prod_{\sigma \in \Gal(L/K)}} L[X] \,:\, f_{\tau\sigma}(1)=\tau(f_{\sigma}(0)), \forall \sigma \in \Gal(L/K) \right\}.
    \]
This means that $X$ is geometrically connected over $K$ and $X_L$ agrees with a chain of affinoid disks $D_\sigma$ indexed by $\sigma \in \Gal(L/K)$ where $0 \in D_{\sigma}$ is glued transversally and over $L$ with $1 \in D_{\tau\sigma}$. The projection $X_L \to X$ is a finite Nisnevich cover that splits itself, hence it is tame relative to $\Spa(K)$, although $L/K$ is not. 
\end{ex}
    

%%%%%%%%%%%%%%%%%%%%%%%%%%%%%%%%%%%%%%%%%%%
\subsection{\texorpdfstring{Alterations and $v$-descent}{Alterations and v-descent}}
\label{ss:alterations-v-descent}

The goal of this subsection is twofold. First, using the good topological properties of qcqs analytic adic spaces (\cref{lem:submersive}), we prove a general descent statement for finite generation of the (tame) fundamental group under surjections (\cref{cor:descent-of-fg-fett}). We then show, using Gabber's uniformisation results, that every rigid space admits a qcqs surjection from a regular rigid space. 

\begin{lem} 
\label{lem:submersive}
    Let $Y\to X$ be a surjective map between qcqs analytic adic spaces. Then $Y\to X$ is topologically submersive, i.e.\ $X$ has the quotient topology. If $Y\to X$ is moreover of finite type, then this holds universally for base-changes along all maps of qcqs analytic adic spaces $X'\to X$. 
\end{lem}

\begin{proof}
As in this setting ``surjectivity is universal'', we are quickly reduced to just proving that $Y\to X$ is topologically submersive. For this, see \cite[Lemma~2.5]{ScholzeECoD} and the remark of \emph{loc.~cit.} that all maps between analytic adic spaces are generalising, see \cite[Lemma~1.1.10(v)]{HuberBook}. 
\end{proof}

Let $Y\to X$ be a map between qcqs rigid spaces. Then pull-back defines a rigid analogue of \labelcref{eq:CechDescentSettingForFEt}, namely a functor 
\[
    \FEt_X\la \Desc(Y/X, \FEt).
\]

\begin{lem} 
\label{lem:univ-top-subm-descent}
    Let $f\colon Y\to X$ be a surjective map between qcqs rigid spaces. Then the pull-back functor
    \[ 
       f^\ast \colon  \FEt_X\la \Desc(Y/X, \FEt)
    \]
    is fully faithful, and the same holds with $\FEtt$ or $\FEtt_{-/S}$ in place of $\FEt$.
\end{lem}

\begin{proof} 
As in \Cref{lem:f-desc-for-FEt-and-FEtt}, if $f^\ast$ is fully faithful for $\FEt$, then $f^\ast$ is also fully faithful for $\FEtt$ and $\FEtt_{-/S}$.
So we have to prove only the assertion for $\FEt$. In this case, the result follows from \cref{lem:submersive} exactly as in \cite[Exp.\ IX, Corollaire 3.3]{SGA1}.
\end{proof}

\begin{rmk}[Effective descent]
In the situation of \cref{lem:univ-top-subm-descent}, if the base field $K$ is of positive or mixed characteristic, one can use the theory of diamonds \cite{ScholzeECoD} to show that    
    \[ 
       f^\ast \colon  \FEt_X\la \Desc(Y/X, \FEt)
    \]
is an equivalence. (We will not use this observation in this paper.)

The argument for this is as follows. In both cases, $X$ is an analytic adic space over $\bZ_p$. By \cite[Lemma 15.6]{ScholzeECoD}, the associated diamond $X^{\diamond}$ is locally spatial and $\FEt_X \simeq \FEt_{X^\diamond}$. Analogous facts hold for $Y$, $Y\times_X Y$, and $Y\times_X Y \times_X Y$ (note that diamantification commutes with fibre products). By Lemma~15.6 of \emph{loc.~cit.}, the map $|Y^\diamond| \to |X^\diamond|$ is still a surjective map of qcqs spaces and by Lemma~12.11 of \emph{loc.~cit.}, this implies that the map of diamonds $f^\diamond \colon Y^\diamond \to X^\diamond$ is surjective, so a v-covering. Any descent datum of diamonds along $f^\diamond$ is effective. By Proposition 10.11 of \emph{loc.~cit.}, the gluing is finite \'etale over $X^\diamond$ and so comes from $\FEt_X$, as desired. We thank Ben Heuer for supplying this argument. 
\end{rmk}

\begin{cor} \label{cor:descent-of-fg-fett}
    Let $Y\to X$ be a surjective map between qcqs rigid spaces over $K$. 
    If for every connected component $Y'$ of $Y$ the fundamental group $\pitame(Y'/K)$ is finitely generated, then for every connected component $X'$ of $X$ the fundamental group $\pitame(X'/K)$ is finitely generated.
\end{cor}

\begin{proof}
Replacing $Y\to X$ with its base change to a connected component of $X$, we may assume that $X=X'$ is connected. \Cref{lem:univ-top-subm-descent} then produces a surjection
\[ 
    \pi_1(\Desc(Y/X), \FEtt_{-/K}) \surj \pitame(X/K).
\]
The source of this map is finitely generated by \cref{prop:DD-fgfp}, and hence so is $\pitame(X/K)$. 
\end{proof}

Rigid spaces are ``$h$-locally regular'' is our next result in \cref{prop:h-locally-regular} below.

\begin{defi}[{Relative analytification, see \cite[Proposition~3.8]{Huber94:Generalization}}] \label{def:relative-analytification}
    Let $R$ be an affinoid $K$-algebra, let $S_0=\Spec(R)$, and let $X_0\to S_0$ be a morphism locally of finite type. Let $S=\Spa(R)$. The {\bf relative analytification} is the morphism of rigid spaces $X=X\times_{S_0} S\to S$. 
\end{defi}

In general, relative analytification preserves open immersions and open coverings (which allows us to globalise the above picture), proper maps \cite[Lemma~5.7.3]{HuberBook}, surjective maps \cite[Lemma~3.9]{Huber94:Generalization} etc.

\begin{prop} 
\label{prop:h-locally-regular}
    Let $X$ be a qcqs rigid space over a perfect non-archimedean field $K$. Then there exists a surjective map of rigid spaces $Y\to X$ with $Y$ qcqs and smooth over $K$.
\end{prop}

This proposition is essentially a result by Berkovich \cite[Theorem 1.3.1]{Berkovich:VanishingCyclesFormalSchemes} but in the adic setting, and in fact it follows from the Berkovich version.
Indeed, any qcqs adic space is taut and, by  \cite[Proposition 8.3.1+Remark 8.3.2]{HuberBook} (cf.\ also \cite[Theorem 1.6.1]{Berkovich:EtaleCohomology}), the category of qcqs rigid spaces over $K$ is equivalent to compact hausdorff strictly $K$-analytic (Berkovich) spaces. Let us denote the quasi-inverse functors realising this equivalence by $(-)^{\mathrm{Berk}}$ and $(-)^{\mathrm{rig}}$. By \cite[Theorem 1.3.1]{Berkovich:VanishingCyclesFormalSchemes}\footnote{The statement of \cite[Theorem 1.3.1]{Berkovich:VanishingCyclesFormalSchemes} gives more information, but rig-smoothness is obtained as a step in the proof.}, there is a ``rig-smooth'' (see \cite[p.~335]{Berkovich:SmoothLocallyContrII}) $Y_0$ with a compact map $Y_0 \to X^\mathrm{Berk}$. So, $Y_0$ is compact and there is a finite set of affinoid subdomains $\{U_i\}$ covering $Y_0$ such that each $U_i$ admits a quasi-\'etale map to $(\mathbb{A}^n_K)^{\rm Berk}$.
By 
\cite[Lemma 1.6.2]{Berkovich:EtaleCohomology} 
and
\cite[Proposition 8.3.1+Remark 8.3.2]{HuberBook} again,
the corresponding rigid space $Y = Y_0^{\mathrm{rig}}$ is qcqs and locally admits an \'etale map to $(\mathbb{A}^n_K)^{\rm an}$. By \cite[Corollary 1.6.10]{HuberBook}, $Y$ is smooth over $K$.

For convenience, let us sketch the central part of the proof of \cref{prop:h-locally-regular} directly in the adic case.

\begin{proof}
Replacing $X$ by the constituents of a finite affinoid open covering, we may assume that $X=\Spa(A)$ is affinoid with $A$ an affinoid $K$-algebra. Let $X_0=\Spec(A)$, which is an excellent scheme. By Gabber's uniformisation theorem \cite[Th\'eor\`eme~1.1]{TravauxGabberIX}, there exists a separated morphism of finite presentation $Y'_0\to X_0$ where $Y'_0$ is regular and which is an $h$-covering. By \cite[Corollary~10.4]{SuslinVoevodsky}, the $h$-covering admits a refinement of normal form, so there exists a proper surjection $X'_0\to X_0$, an open covering $X'_0 = \bigcup_{i=1}^r U_{0,i}$, and a commutative square
\[ 
    \begin{tikzcd}
        \coprod U_{0,i} \ar[d] \ar[r] & Y'_0 \ar[d] \\
        X'_0 \ar[r] & X_0.
    \end{tikzcd}
\]
Analytifying the above square in a relative sense, we obtain a commutative square of rigid spaces 
\[ 
    \begin{tikzcd}
        \coprod U_i \ar[d] \ar[r] & Y' \ar[d] \\
        X'\ar[r] & X
    \end{tikzcd}
\]
where $X'\to X$ is a proper  (cf.~\cite[Lemma~5.7.3]{HuberBook}) surjection and $\coprod U_i \to X'$ is a Zariski open covering, and where $Y'$ is regular (by which we mean that the local rings at classical points are regular).
Indeed, one checks (using e.g.\ \cite[Satz 2.1]{Kopf}) that relative analytification preserves completed local rings of classical points. Cf.\ also \cite[Lemma 2.6.3]{Berkovich:EtaleCohomology} for the Berkovich spaces incarnation of this result and \cite[I \S 4.1, Proposition 2]{Bosch} for a proof of the same statement when the target is $\Spa(K)$. Then one checks that regularity together with perfectness of $K$ implies smoothness (cf.\ \cite[\S 5]{Berkovich:VanishingCyclesAnalytic}).

Since $X'\to X$ is proper, the rigid space $X'$ is quasi-compact and separated. Thus, there exist quasi-compact opens $V_i\subseteq U_i$ such that $X'=\bigcup_{i=1}^r V_i$. To see this, take a refinement of $\coprod U_i \to X'$ by affinoids, hence quasi-compact opens, and let $V_i$ be the union of those opens inside $U_i$. The image of $\coprod _{i=1}^r V_i\to Y'$ is quasi-compact, and hence contained in a quasi-compact open $Y\subseteq Y'$, which is also quasi-separated because $Y'\to X$ is separated and $X$ is affinoid. The map $Y\to X$ is surjective by construction. 
\end{proof}

The following example explains a subtle point in the above proof. 

\begin{ex}[Punctured blowup]
Let $X = \Spa(K\langle x, y\rangle)$ be the affinoid bidisc, and let $Y'\to X$ be the blowup of the ideal $(x, y)$. Let $P \in Y'$ be a $K$-point in the exceptional divisor, and let $Y = Y'\setminus\{P\}$. Then the map $Y\to X$ is surjective, but no quasi-compact open $U\subseteq Y$ surjects onto $X$. Suppose otherwise, and consider in suitable coordinates the chart $\Spa(K\langle x,t\rangle)\subseteq Y'$ with $t=y/x$ in which $p$ has coordinates $(0,0)$. Then for $n\gg 0$, the open ball $B_n = \{|x|, |t|<1/n\}\subseteq Y'$ is disjoint from $U$, and therefore the points $(x,0) \in X$ with $0<|x|<1/n$ are not in the image of $U\to X$.  

Note that if $Y'_0\to X_0=\Spec(K\langle x, y\rangle)$ is the blowup of $(x,y)$, and $Y_0 = Y'_0\setminus\{P\}$, then $Y\to X$ is the analytification of $Y_0\to X_0$. However, $Y_0\to X_0$ is not an $h$-covering even though it is surjective. The refinement of normal form used in the proof of \cref{prop:h-locally-regular} is needed precisely to ensure that after taking analytifications, the map will remain surjective after passing to a suitable quasi-compact open.
\end{ex}

%%%%%%%%%%%%%%%%%%%%%%%%%%%%%%%%%%%%%%%%%%%
\subsection{Review of log formal schemes}
\label{ss:log-formal-schemes}

In this subsection, we discuss the log geometry of formal schemes needed for the treatment of semistable formal models of rigid spaces in \S\ref{ss:log-abhyankar}. See \cite[\S 7]{Fersi} for some foundations of log formal schemes (though in lesser generality than needed here).

As before, we fix a complete rank one valuation ring $K^\circ$ with fraction field $K$ and pick a pseudouniformiser $\pi\in K^\circ$. All formal schemes over $K^\circ$ are assumed to be $\pi$-adic, i.e.\ such that $\pi$ generates an ideal of definition. Logarithmic structures can be defined on any ringed topos, and by a {\bf log formal scheme} we mean a formal scheme $\mathfrak{X}$ endowed with a log structure 
\[
\alpha\colon\cM_\fX\to\cO_\fX.
\]

The data of a $\pi$-adic formal scheme $\fX$ is equivalent to the data of a sequence of schemes $\fX_n = \fX\otimes K^\circ/\pi^{n+1}$ over $K^\circ/\pi^{n+1}$ together with identifications
\[ 
    \fX_n \simeq \fX_{n+1} \otimes_{K^\circ/\pi^{n+2}} K^\circ/\pi^{n+1}.
\]
The following basic result extends this equivalence to log structures.

\begin{lem} \label{lem:log-fsch-equiv}
    Let $\cM_\fX$ be a log structure on a $\pi$-adic formal scheme $\fX$, and let $\cM_{\fX_n}$ be the induced log structure on the scheme $\fX_n$. Then
    \[
    \cM_{\fX_n} = \cM_\fX/U_n
    \quad \text{where} \quad
    U_n = \ker(\cO_{\fX}^\times \to \cO_{\fX_n}^\times) = 1+\pi^{n+1}\cO_\fX
    \]
    and
    \[ 
        \cM_\fX = \varprojlim \cM_{\fX_n} = \varprojlim \cM_\fX/U_n.
    \]
    Conversely, given a sequence of log structures $\cM_{\fX_n}$ on $\fX_n$ and identifications $\cM_{\fX_{n+1}}|_{\fX_{n}} \simeq \cM_{\fX_n}$, the above formula defines a log structure on $\fX$ whose restriction to $\fX_n$ equals $\cM_{\fX_n}$. 
\end{lem}

\begin{proof}
The preimage of $\cO_{\fX_n}^\times\subseteq\cO_{\fX_n}$ under $\alpha_n\colon \cM_{\fX}\to\cO_\fX\to\cO_{\fX_n}$ equals $\cO_\fX^\times$. Therefore
\[
    \cM_{\fX_n} = \cM_\fX \oplus_{\alpha_n^{-1}(\cO_{\fX_n}^\times)} \cO_{\fX_n}^\times = \cM_{\fX} \oplus_{\cO_{\fX}^\times} (\cO_{\fX}^\times/U_n) = \cM_\fX/U_n,
\]
showing the first assertion. The second follows from the proof of \cite[Proposition~7.8]{Fersi}. 
We omit the proof of the final assertion.
\end{proof}

We are primarily interested in log structures which locally admit a chart, in particular smooth log formal schemes over $K^\circ$. For a monoid $P\to A$ mapping into a $\pi$-adic ring $A$, we denote by $\Spf(P\to A)$ the log scheme $\Spf(A)$ endowed with the log structure charted by $P$. 

\begin{defi}
    \begin{enumerate}[(a)]
        \item A log formal scheme $\fX$ is {\bf quasi-coherent} (resp.\ {\bf integral}, {\bf saturated}, {\bf fine}, {\bf fs}) if it is \'etale locally of the form $\Spf(P\to A)$ for a monoid (resp.\ integral, saturated, fine, fs monoid) $P$ mapping into a $\pi$-adic ring $A$.
        \item (See \cite[Definition~3.5.1]{AHLS-Log} for the scheme case.) Let $f\colon \fY\to\fX$ be a morphism of saturated adic log formal schemes over $K^\circ$. We say that $f$ is {\bf smooth} (resp.\ {\bf \'etale}, {\bf Kummer \'etale}) if for every strict \'etale map $\fX' = \Spf(P\to A)\to\fX$ with a saturated monoid $P$, \'etale locally on $\fY' = \fY\times_\fX\fX'$ there exists a  smooth (resp.\ \'etale, resp.\ Kummer \'etale) morphism of monoids $P\to Q$ (with $\Sigma$ the set of primes non-invertible on $X$) such that $\fY'\to\fX'$ factors as
    \[
        \fY\la \Spf\left(Q\to (A\otimes_{\bZ[P]}\bZ[Q])^\wedge_\pi\right) \to \Spf(P\to A) = \fX'
    \]
    where the first map is strict and smooth (resp.\ \'etale, resp. \'etale). Here, $(-)^\wedge_\pi$ denotes the $\pi$-adic completion.  
    \end{enumerate}
\end{defi}

Arguing as in the scheme case \cite[\S 3.5]{AHLS-Log}, one shows these three classes of morphisms (smooth, \'etale and Kummer \'etale) are closed under composition and base change. We warn the reader, however, that, since smooth maps of monoids are not always finitely generated,  the underlying maps of formal schemes might not be locally topologically of finite type in general. Below we shall work exclusively with Kummer \'etale maps between smooth log formal schemes over $\Spf(K^\circ)^{\rm log}$ where $K$ is algebraically closed or discretely valued, in which case there is no problematic behaviour. 

\begin{defi} \label{def:std-log-str}
    The log structure ``induced by the generic fibre'' is the key example of a log structure on a formal scheme.
    Let $\fX$ be a $\pi$-adic formal scheme over $K^\circ$. 
    \begin{enumerate}
        \item The subsheaf
        \[ 
            \cM_\fX^{\rm std} = \cO_\fX \times_{\cO_\fX[1/\pi]} \cO_\fX[1/\pi]^\times \quad\subseteq\quad \cO_\fX
        \]
        is a log structure, called the {\bf standard log structure}.
        \item We denote by $\fX^{\rm log}$ the log formal scheme $(\fX, \cM_\fX^{\rm std})$, and by $\fX^{\rm log}_k$ its special fibre $\fX_k$ with the induced log structure $\cM_\fX^{\rm std}|_{\fX_k}$. 
        \item A log structure $\cM$ on $\fX$ is {\bf vertical} if the image of $\cM\to\cO_\fX$ lies in $\cM_\fX^{\rm std}$, i.e.\ if the image of every local section of $\cM$ in $\cO_\fX[1/\pi]$ is invertible.
    \end{enumerate}
\end{defi}

Intuitively, being vertical means that the induced log structure on the rigid generic fibre $\fX_{\rm rig}$ is trivial.  Thus, $\cM_\fX^{\rm std}$ is the maximal vertical log structure. An important special case is the log formal scheme $\Spf(K^\circ)^{\rm log}$, which admits a chart by the monoid 
    \[
        V = K^\circ \cap K^\times \la K^\circ,
    \]
    cf.\ \cite[Lemma~4.4.1]{AHLS-Log}. Note that $V$ is a valuative monoid, which is moreover divisible if $K$ is algebraically closed. Thus $\Spf(K^\circ)_k^{\rm log} = \Spec(V\to k)$ is of type $\typeVd$ if $K$ is algebraically closed or if $K$ is discretely valued.

\begin{lem} \label{lem:spf-Kcirc-satisfies364}
    The valuation map $\nu \colon K^0 \cap K^\times \to \Gamma_K^+$ admits a splitting $\sigma$, which induces an isomorphism of log formal schemes
    \[
    \Spf(K^\circ)^{\log} \isomto \Spf(\Gamma_K^+ \xrightarrow{\sigma} K^\circ).
    \]
\end{lem}
\begin{proof} 
    The exact sequence 
    \[
        \begin{tikzcd}
            1 \ar[r] & K^{\circ\times} \ar[r] & K^\times \ar[r,"\nu"] & \Gamma_K\ar[r] & 0.
        \end{tikzcd}
    \]
   splits because by assumption $\Gamma_K \simeq \bZ$ is free or $K^{\circ \times}$ is divisible. The restriction of a splitting induces a splitting $\sigma \colon \Gamma_K^+ \to K^\circ \cap K^\times$. 
\end{proof}


The most important special case of a log smooth formal scheme over $\Spf(K^\circ)^{\rm log}$ is $\fX^{\rm log}$ for a semistable formal scheme $\fX$ over $K^\circ$. Before we show this, let us recall the definition. 

\begin{defi} \label{def:ss-formal-sch}
    A formal scheme $\fX$ over $K^\circ$ is {\bf semistable} if \'etale locally on $\fX$ there exists a pseudouniformiser $\pi\in K^\circ$ and an \'etale map
    \[
        \fX \la \Spf\left( K^\circ\langle T_1,\dots, T_n\rangle/(\pi - T_1\cdot\ldots\cdot T_n) \right).
    \]
    We say that $\fX$ is {\bf strictly semistable} if such a map exists Zariski locally on $\fX$. 
\end{defi}

\begin{ex} \label{ex:standard-log-structure-Kcirc}
    Here is an example showing that our definition of a strictly semistable formal scheme covers the usual cases. For $n \geq s > 0$, the morphism
    \[
    \Spf(K^\circ\langle X_1,\ldots, X_s, X_{s+1}^\pm, \ldots, X_n^\pm\rangle/(\pi - X_1 \cdot \ldots \cdot X_s)) \to \Spf( K^\circ\langle T_1,\dots, T_n\rangle/(\pi - T_1\cdot\ldots\cdot T_n))
    \]
    sending $X_i \mapsto T_i$ for $i \not= s$ and $X_s \mapsto T_s \cdot(T_{s+1} \cdot \ldots \cdot T_n)$ is an open immersion. In particular, its source is strictly semistable. 

    Note that the case $s=1$ covers the formally smooth (good reduction) case as well.
\end{ex}

We define the standard semistable monoid in $n$ variables relative to a monoid $M$ and $\gamma \in M$ (cf.\ \cite[Example~2.4.5]{AHLS-Log}) as the monoid $P_n(\gamma)$ together with a smooth map 
\[
M \la P_n(\gamma) = M[e_1,\ldots,e_n]/(e_1 + \ldots + e_n = \gamma).
\]

\begin{lem} \label{lem:strictly-semistable-formal-scheme-reduces-to-setup364}
Let $M = \Gamma_K^+$ and choose a splitting $\sigma$ as in \cref{lem:spf-Kcirc-satisfies364}. 
Let $\fX$ be a (strictly) semistable formal scheme over $K^\circ$. Then the standard log structure on $\fX$ admits \'etale  locally (resp.\ Zariski locally), with a pseudouniformiser $\pi$ that satisfies 
$\sigma(\nu(\pi)) = \pi$, charts of the form 
\[
    \begin{tikzcd}
        \fX^{\log} \arrow[rr,"\text{strict \'etale}"] \arrow[d]
        && \Spf\left(P_n(\nu(\pi)) \xrightarrow{\alpha} K^\circ\langle T_1,\dots, T_n\rangle/(\pi - T_1\cdot\ldots\cdot T_n) \right) \arrow[d] \\
        \Spf(K^\circ)^{\log}  \arrow[rr,"\sim"] & &  
        \Spf(M \xrightarrow{\sigma} K^\circ) 
    \end{tikzcd}
\]
where $\alpha(e_i) = T_i$, and $\alpha$ agrees with $\sigma$ on $M = \Gamma_K^+$.
Consequently, in the strictly semistable case the special fibre $\fX^{\log}_k$ is described by \cref{setup:surgery}. 
\end{lem}

\begin{proof}
    By \cref{def:ss-formal-sch} it suffices to understand a global chart for the standard log structure on $\fX = \Spf\left( K^\circ\langle T_1,\dots, T_n\rangle/(\pi - T_1\cdot\ldots\cdot T_n) \right)$. For the chosen splitting $\sigma$ the pseudouniformiser $\pi$ in the equation defining $\fX$ might differ from $\sigma(\nu(\pi))$ by a unit in $K^\circ$. After scaling one of the coordinates appropriately, we may assume that indeed $\sigma(\nu(\pi)) = \pi$. The rest is obvious.     
\end{proof}

Let us recall Temkin's uniformisation result, to be used in the proof of \cref{ss:tamepi1-rig-fgfp}.

\begin{thm}[{Temkin \cite[Corollary~3.3.2]{Temkin2017:AlteredLocalUniformization}}] \label{thm:Temkin-unif}
    Let $X$ be a smooth qcqs rigid space over a non-archimedean field $K$. Then there exists a finite extension $L$ of $K$, an affine strictly semistable formal scheme $\fY$ over $L^\circ$, and a surjective \'etale map of rigid $L$-spaces $\fY_\rig\to X_L$.
\end{thm}

The following proposition lists some good properties of certain smooth log formal schemes over $\Spf(K^\circ)^{\rm log}$.

\begin{prop} \label{prop:log-sm-vertical}
    Suppose that $K$ is algebraically closed or discretely valued. Let $\fX$ be a log formal scheme over $\Spf(K^\circ)^{\rm log}$.
    \begin{enumerate}[(a)]
        \item \label{propitem:log-sm-vertical-a}
        If $\fX$ is smooth over $\Spf(K^\circ)^{\rm log}$, then its underlying formal scheme is normal and flat over $K^\circ$, in particular $\eta$-normal.
        
        \item \label{propitem:log-sm-vertical-b}
        If $\fX$ is smooth and vertical over $\Spf(K^\circ)^{\rm log}$, then $\cM_\fX$ coincides with the standard log structure $\cM^{\rm std}_\fX$.
        
        \item \label{propitem:log-sm-vertical-c}
        If the underlying formal scheme of $\fX$ is semistable over $K^\circ$ and $\cM_\fX$ is the standard log structure, then $\fX$ is smooth and vertical over $\Spf(K^\circ)^{\rm log}$. 
        
        \item \label{propitem:log-sm-vertical-d}
        If $\fX$ is smooth and vertical over $\Spf(K^\circ)^{\rm log}$ and $\fY\to \fX$ is Kummer \'etale, then $\fY$ is smooth and vertical over $\Spf(K^\circ)^{\rm log}$.
    \end{enumerate}
\end{prop}

\begin{rmk}
We expect that if $K$ is algebraically closed and $\fX$ is log smooth and vertical over $\Spf(K^\circ)^{\rm log}$, then there exists a log blow-up $\fX'\to\fX$ such that the underlying scheme of $\fX'$ is semistable and $\cM_{\fX'}$ is the standard log structure.
\end{rmk}

\begin{proof}
For log schemes over $K^\circ$, these assertions are proved in \cite[\S5]{AHLS-Log}. Here we show how to deduce from them the log formal scheme versions.

\ref{propitem:log-sm-vertical-a}
Working locally, we may assume that $\fX$ is the $\pi$-adic completion of an affine charted smooth log scheme $Y = \Spec(B)$ over $\Spec(V\to K^\circ)$, so that $\fX = \Spf(\widehat{B})$ where $\widehat{B}$ is the $\pi$-adic completion of $B$. By \cite[Lemma 5.3.4]{AHLS-Log}, log smoothness implies that the ring $B$ is normal and flat over $K^\circ$. 
By \cite[Lemma B.6 + Remark B.7]{GabberZavyalov}, the map $B \to \widehat{B}$ is ind-smooth. In particular, it is flat. Moreover, normality persists under smooth maps (see \stacks[Lemma]{033C}) and cofiltered limits (see \stacks[Lemma]{037D}), so we get that $\widehat{B}$ is normal and flat over $K^\circ$, as desired.


\ref{propitem:log-sm-vertical-b} Working locally we may assume that there exists a smooth homomorphism $M\to P$ and a strict \'etale map $X\to \Spf(P\to \widehat{B})$ where $\widehat{B}$ is the $\pi$-adic completion of $B = K^\circ\otimes_{\bZ[M]}\bZ[P]$. We may replace $X$ with the target of this map, in which case $X$ is the $\pi$-adic completion of the smooth log scheme $X_0 = \Spec(P\to B)$. Note that under our assumptions on $K$, the underlying scheme of $X_0$ is of finite presentation over $K^\circ$. 
Moreover, for every strict \'etale $U\to X$ there exists a strict \'etale $U_0\to X_0$ whose $\pi$-adic completion is $U\to X$ (see \cite[1.7.3]{HuberBook}). This enables us to work \'etale locally on $X$. 

By the scheme-theoretic result \cite[Proposition~5.3.1]{AHLS-Log}, we know that every $f\in B$ which satisfies $fg = \pi^n$ is locally of the form $\alpha(p)u$ for $p\in P$ and $u\in B^\times$, where $\alpha\colon P\to B$ is the given map. We want to prove the analogous result for $\widehat{B}$. Let $f\in \widehat{B}$ satisfy $fg = \pi^n$ for some $g\in \widehat{B}$ and $n\geq 0$ and let $J = f\widehat{B}$. Since $\pi^n\in J$ and $\widehat{B}/\pi^n = B/\pi^n$, there exists an ideal $J_0\subseteq B$ such that $J_0\cdot \widehat{B} = J$. Since $B\to B[1/\pi]\times\widehat{B}$ is faithfully flat (by  \cite[Chapter~0, Proposition~8.2.17]{FujiwaraKato}), and since $J_0\cdot B[1/\pi] = B[1/\pi]$ and $J_0\cdot \widehat{B} = J = (f)$ are both invertible ideals, by flat descent for line bundles we see that $J_0$ is an invertible ideal. Localising on $X$ and $Y$, we may assume that $J_0 = (f_0)$ is principal. Since $J_0$ contains $\pi^n$, we can write $\pi^n = f_0g_0$, and hence (shrinking $X$ further) $f_0 = \alpha(p)u_0$ for some $p\in P$ and $u_0\in B^\times$. Since $f \widehat{B} = f_0\widehat{B}$, there exists a unit $u_1\in \widehat{B}^\times$ such that $f=u_1 f_0$. Then $f = \alpha(p)u$ where $u=u_0u_1\in \widehat{B}^\times$. 

\ref{propitem:log-sm-vertical-c} This can be deduced directly from the scheme version, \cite[Corollary~5.3.7]{AHLS-Log}.

\ref{propitem:log-sm-vertical-d} Since Kummer \'etale maps are smooth and smooth maps are closed under composition, we have that $\fY$ is smooth. It remains to show $\fY$ is vertical. We may assume that $\fY = \fX\times_{\Spf(P\to \bZ[P])}\Spf(Q\to\bZ[Q])$ for a Kummer \'etale map $P\to Q$. We must show that the image $g$ of every $q\in Q$ in $\cO_\fY[1/\pi]$ is invertible. Since $P\to Q$ is Kummer \'etale, it is injective and we have $nq=p\in P$ for some $n\geq 1$ invertible in $K^\circ$. Since $\fX$ is vertical, the image $f$ of $p$ in $\cO_\fX[1/\pi]$ is invertible. But $g^n=f$, and it follows that $g$ is invertible as well.
\end{proof}

\begin{prop} \label{prop:FEt-formal-vs-sp-fibre}
    Let $\fX$ be a saturated log formal scheme over $K^\circ$. Then the inclusion $\fX_k\to \fX$ induces an equivalence
    \[
        \FEt_\fX \isomto \FEt_{\fX_k}.
    \]
\end{prop}

\begin{proof}
Let $\fX_n$ be the reduction of $\fX$ modulo $\pi^{n+1}$. It is clear from the definition that if $\fY\to\fX$ is finite Kummer \'etale, then so are $\fY_n\to\fX_n$ and $\fY_k\to \fX_k$, so that our functor is well-defined. By \cref{prop:FKEt-top-inv}, the restriction functors $\FEt_{\fX_n}\to\FEt_{\fX_k}$ are equivalences. On the other hand, by \cref{lem:log-fsch-equiv}, giving a map of formal log schemes $\fY\to\fX$ is the same as giving a compatible family of maps $\fY_n\to \fX_n$. This implies that the functor is fully faithful, and hence it remains to prove that it is essentially surjective \'etale locally on $\fX$. We may thus assume that $\fX=\Spf(P\to A)$ for a saturated monoid $P$. Let $Y\to \fX_k$ be an object of $\FEt_{\fX_k}$. By the local structure in \cref{lem:fKet-local-structure} we may assume that there exists a finite collection of Kummer \'etale monoid maps $P\to Q_j$, $j=1,\dots, r$, such that $Y = \coprod_{j=1}^r \Spec(Q_j \to (A\otimes k)\otimes_{\bZ[P]}\bZ[Q_j])$. We may then define 
\[
    \fY = \coprod_{j=1}^r \Spf(Q_j \to (A\otimes_{\bZ[P]}\bZ[Q_j])^\wedge_\pi) \la \Spf(P\to A) = \fX
\]
which is an object of $\FEt_\fX$ whose restriction to $\fX_k$ equals $Y\to\fX_k$.
\end{proof}


%%%%%%%%%%%%%%%%%%%%%%%%%%%%%%%%%%%%%%%%%%%
\subsection{Tame covers and Kummer covers (I): purity}
\label{ss:log-abhyankar}

In this subsection, $K$ will be either algebraically closed  or 
    discretely valued.
Our goal is to prove that for a semistable formal scheme $\fX$ over $K^\circ$, a finite \'etale cover $Y\to X$ of its generic fibre $X=\fX_{\rig}$ is tame (in the absolute sense, i.e.\ relative to $X$) if and only if it extends to a Kummer \'etale cover of $\fX^{\rm log}$ ($\fX$ endowed with the standard log structure). In order to do this, we shall first establish a version for schemes over~$K^\circ$, which in turn will ultimately rely on the classical log purity (a.k.a.\ log Abhyankar's lemma) recalled below. This is particularly delicate in the non-noetherian setting (i.e.\ if $K$ is algebraically closed), as we need to perform a rather intricate approximation argument to be able to apply classical log purity. Let us mention here that the case of $\Spec(K^+)$ for an arbitrary valuation ring $K^+$ has been handled in \cite[Corollary~4.4.9]{AHLS-Log}, identifying the tame Galois group of $K$ with the Kummer \'etale fundamental group of $\Spec((K^+\setminus 0)\to K^+)$. 

To state log purity, we need the following terminology. If $U\to X$ is an \'etale morphism of schemes and $\zeta\in X$ is a point such that $\cO_{X,\zeta}$ is a valuation ring, we say that $U\to X$ is {\bf tame above $\zeta$} if the map $\Spa(U, X)\to \Spa(X, X)$ is tame at every $u\in \Spa(U, X)$ mapping to the point $x\in \Spa(X, X)$ corresponding to the valuation induced by $\zeta$.
Explicitly, we are testing tameness of $\Spa(U,X) \to \Spa(X,X)$ at all preimages of the root triple $(\eta, \cO_{X,\zeta},\zeta)$ of $\Spa(X,X)$ where $\eta\in X$ is the image of the generic point of $\cO_{X,\zeta}$.

A natural setting where tameness above~$\zeta$ comes up is when $X$ is the underlying scheme of a regular log scheme and $U \to X$ factors through a finite étale morphism $U \to X^*$ (the locus of trivial log structure).
Since~$X$ is normal by \cite[Theorem~4.1]{KatoToricSingularities}, the local rings at the generic points of $X \setminus X^*$ are valuation rings and it makes sense to ask whether $U \to X$ is tame above these points.

\begin{thm}[{Classical log purity, \cite[Theorem~13.3.43]{GabberRameroFoundations}, \cite[Theorem~B]{Mochizuki}}] 
\label{thm:GR-log-Abhyankar}
    Let $X$ be a regular log scheme, and 
    let $X^* \subseteq X$ be the open subset where the log structure is trivial. 
    
    Then the restriction functor $\FEt_X\to \FEt_{X^*}$ induces an equivalence with the full subcategory of $\FEt_{X^*}$ consisting of \'etale covers that are tame above every generic point of $X\setminus X^*$.
\end{thm}

Our first result is an extension of the above theorem to semistable log schemes over $K^\circ$, i.e.\ those which \'etale locally admit an \'etale morphism to $\Spec(K^\circ[T_1,\dots,T_n]/(\pi- T_1\cdot\ldots\cdot T_s))$ for some $n \geq s \geq 0$ and some pseudouniformiser $\pi\in K^\circ$. To this end, we first need to check that the local rings of the generic points of the special fibre are valuation rings. This follows from the lemma below, in which we use the following terminology analogous to \cref{defi:eta-normal}: a scheme $X$ locally of finite type over a valuation ring $K^+$ is {\bf $\eta$-normal} if it is locally of the form $\Spec(A)$ where $A$ is a finitely generated and flat $K^+$-algebra which is integrally closed in $A\otimes_{K^+} K$ where $K$ is the field of fractions of $K^+$. We note that being $\eta$-normal is an \'etale local property thanks to \stacks[Lemma]{03GG}.


\begin{lem} \label{lem:local-ring-val-ring}
    Let $(K,K^+)$ be a valued field with residue field $k = K^\succ$ and pseudouniformiser~$\pi$.
    Let $X$ be a reduced scheme that is of finite type, flat and $\eta$-normal over~$K^+$.
    Then for every generic point~$\beta$ of the special fibre~$X_k$, the local ring~$\cO_{X,\beta}$.
    Moreover, the map $K^+\to \cO_{X,\beta}$ is an extension of valuation rings with finite ramification index (see \stacks[Definition]{0ASG}). 
\end{lem}

\begin{proof}
We discuss two special cases and deduce from them the general statement.
First, we treat the case where $X = \bA^n_{K^+} = \Spec (K^+[T_1,\ldots,T_n])$ and $\beta$ is the unique generic point of the special fibre corresponding to the prime ideal $\mathfrak{m}_{K^+}\cdot K^+[T_1,\ldots,T_n]$.
The field of fractions of $\cO_{X,\beta}$ is $K(T_1, \dots, T_n)$, and a fraction $f/g$ with $f,g\in K^+[T_1,\dots, T_n]$ belongs to $\cO_{X,\beta}$ if and only if we can find a representation with $g\notin \mathfrak{m}_{K^+} \cdot K^+[X_1, \ldots, X_n]$.
Since every element of $K^+[T_1,\ldots,T_n]$ can be written as a product $x h$ where $x\in K^+$ and $h\notin \mathfrak{m}_{K^+}[T_1,\ldots,T_n]$ (thus $h\in \cO_{X,\beta}^\times$), it follows that $\cO_{X,\beta}$ is a valuation ring and that  $K^+\to \cO_{X,\beta}$ is an extension of valuation rings with valuation index one. 

Second, we treat the case where~$K^+$ is henselian and~$X$ is finite over~$K^+$.
Then $\cO_{X,\beta}$ is reduced and finite flat over~$K^+$.
In this situation $\eta$-normality implies normality, so $\cO_{X,\beta}$ is normal and this already implies that it is a valuation ring.
By the finiteness of~$\cO_{X,\beta}$ over~$K^+$, it follows that $K^+\to \cO_{X,\beta}$ is an extension of valuation rings with finite valuation index.

For the general case, since the assertion is local on $X$, we may assume that $X = \Spec(A)$ is affine and $X_k$ is irreducible, so $\beta$ is the unique generic point of $X_k$.
By the Noether normalisation lemma, there exist $f_1,\dots, f_n\in A$ such that the map $f=(f_1,\dots, f_n)\colon X\to \bA^n_{K^+}$ induces a finite map $f_k\colon X_k\to \bA^n_k$ which sends $\beta$ to the generic point $\gamma$.
Let
\[
    f_\beta \colon X_\beta^h \la \bA_{K^+,\gamma}^{n,h}
\]
be the induced map of henselisations.
In the first paragraph we have seen that the localisation $\bA_{K^+,\gamma}^n$ -- and hence also its henselisation -- is the spectrum of a valuation ring extending $K^+$ with finite ramification index.
The second paragraph deduces that the same holds for~$X_\beta^h$.
Since a local ring is a valuation ring if and only if its henselisation is, we conclude that $\cO_{X,\beta}$ is a valuation ring extending $K^+$ with finite ramification index.
\end{proof}

Henceforth in this section we use the following notation. If $Z$ is a scheme over $K^\circ$, we denote by $Z^*$ its generic fibre $Z\otimes_{K^\circ} K$. We use the natural log scheme variants of \cref{def:std-log-str}.  Our version of \cref{thm:GR-log-Abhyankar} over $K^\circ$ is the following.    

\begin{thm}[Purity, scheme version] \label{thm:scheme-Abhyankar}
  Let $X$ be a semistable scheme over $K^\circ$ and let $Y^*\to X^* = X_K$ be a finite \'etale map.  Assume that $K$ is either algebraically closed or discretely valued. 
 Let $Y\to X$ be the normalisation of $X$ in $Y^*$. The following are equivalent:
  \begin{enumerate}[(a)]
    \item \label{propitem:Abhy-tame} 
        $Y^*\to X^*$ is tame relative to $X$,
    \item \label{propitem:Abhy-tame-boundary}
        $Y^*\to X^*$ is tame above every generic point of the special fibre $X_k$,
    \item \label{propitem:Abhy-ket}
        the morphism $Y^{\rm log}\to X^{\rm log}$ is Kummer \'etale.
  \end{enumerate}
\end{thm}

\noindent\emph{Proof (part one).} 
\label{proof:scheme-Abhyankar}
Assertion \labelcref{propitem:Abhy-tame} obviously implies assertion \labelcref{propitem:Abhy-tame-boundary}. Assume \labelcref{propitem:Abhy-ket} and show that it implies \labelcref{propitem:Abhy-tame}. For any $x \in \Spa(X^*,X)$ we endow $S = \Spec(k(x)^+)$ with the standard log structure. The resulting  map of log schemes $\tilde{x} = \Spec(k(x)^+ \setminus 0 \to k(x)^+) \to X$ yields a pull-back of $Y \to X$ to $\tilde{x}$ that is a Kummer \'etale cover of $\tilde{x}$.
By  \cite[Corollary 4.4.9]{AHLS-Log} this corresponds to a tamely ramified extension of $k(x)$, and this shows \labelcref{propitem:Abhy-tame}.

For \labelcref{propitem:Abhy-tame-boundary}$\Rightarrow$\labelcref{propitem:Abhy-ket}, suppose first that $K$ is discretely valued. In this case, $\Spec(K^\circ)$ is log regular for the standard log structure. By \cite[Corollary~5.3.7]{AHLS-Log} (the scheme version of  \cref{prop:log-sm-vertical}(c)) the map $X\to \Spec(K^\circ)$ is log smooth for the standard log structures on source and target. Therefore by \cite[Theorem~8.2]{KatoToricSingularities} the log scheme $X$ is log regular as well. By log purity (\cref{thm:GR-log-Abhyankar}) $Y^*\to X^*$ extends to a Kummer \'etale $Y'\to X$. It remains to check that $Y'=Y$, for which we note that the log structure on $Y'$ is also the standard one, and again $Y'$ is log regular (being smooth over $X$). Thus in particular $Y'$ is normal \cite[Theorem~4.1]{KatoToricSingularities}, and hence $Y=Y'$.

It remains to show \labelcref{propitem:Abhy-tame-boundary}$\Rightarrow$\labelcref{propitem:Abhy-ket} for $K$ algebraically closed. In this case, log purity is not readily available, but we shall ultimately be able to apply it using approximation techniques. 

\hphantom{a} \hfill \emph{...the proof will resume after \cref{lem:generic-points-ugggh}. $\square$}  

\medskip

The approximation part of the proof of \cref{thm:scheme-Abhyankar} will rely on the following lemma. 

\begin{lem}[Zavyalov] \label{lem:Zavyalov}
    Let $K$ be an algebraically closed non-archimedean field, and let $\pi\in K$ be a pseudouniformiser. Then $K^\circ$ is isomorphic to the filtered colimit $\varinjlim R_\lambda$ of subalgebras $R_\lambda\subseteq K^\circ$ such that
    \begin{enumerate}[(i)]
        \item $\pi \in R_\lambda$;
        \item $R_\lambda$ is a strictly henselian local ring and $R_\lambda\to K^\circ$ is a local homomorphism;
        \item $R_\lambda$ is excellent and regular;
        \item the divisor $V(\pi)_{\rm red}\subseteq \Spec(R_\lambda)$ has strict simple normal crossings.
    \end{enumerate}
\end{lem}

\begin{proof}
We apply \cite[Lemma~A.2]{Zavyalov} and replace the resulting subrings of $K^\circ$ with their respective strict henselisations with respect to the map to the residue field $k$.
\end{proof}

Let us henceforth fix a pseudouniformiser $\pi$ of $K^\circ$. We will consider schemes $Z$ defined over subalgebras $R\subseteq K^\circ$ containing $\pi$. In this case, we shall extend our previous notation and denote by $Z^*$ the open subset $D(\pi)\subseteq Z$ (this agrees with the previous use since $K= K^\circ[1/\pi]$).

\begin{lem}  
\label{lem:straightforward-approx}
    Let $K$ be an algebraically closed non-archimedean field. Let $X=\Spec(A)$ be an affine $K^\circ$-scheme equipped with an \'etale map
    \[ 
        X\longrightarrow W = \Spec(K^\circ[T_1, \dots, T_n]/(\pi- T_1\cdot\ldots\cdot T_s)).
    \]
    Let $X^* = X_K = \Spec(A[1/\pi])$ be its generic fibre and let $Y^*=\Spec(B)\to X^*$ be a finite \'etale map. Let $\{R_\lambda; \lambda \in I\}$ be a cofiltering system of subalgebras of $K^\circ$ containing $\pi$ such that $K^\circ=\varinjlim R_\lambda$, 
    and let
    \[
        W_\lambda = \Spec(R_\lambda[T_1, \dots, T_n]/(\pi- T_1\cdot\ldots\cdot T_s)).
    \]
    Then for some $\lambda\in I$ there exist an \'etale map $X_\lambda=\Spec(A_\lambda)\to W_\lambda$ and a finite \'etale map $Y_\lambda=\Spec(B_\lambda)\to \Spec(A_\lambda[1/\pi]) = X^*_\lambda$ whose base changes to $K^\circ$ are $X\to W$ and $Y^*\to X^*$, respectively.
\end{lem}

\begin{proof}
Since \'etale morphisms of affine schemes are of finite presentation, this follows from \cite[8.8.2 and 8.10.5(viii,x)]{EGAIV3} and \stacks[Lemma]{07RP}.
\end{proof}

We now fix a cofiltering system $\{R_\lambda\}$ as in \cref{lem:Zavyalov} and assume that there is an initial index $\lambda_0$ for which the claim of \cref{lem:straightforward-approx} holds. Then we define $X_\lambda = \Spec(A_\lambda) \to W_\lambda$ and $Y_\lambda = \Spec(B_\lambda) \to X_\lambda^*$ by base change from the objects for the initial $\lambda_0$. 

\begin{lem}
\label{lem:locus of tameness}
     Assume that $K$ is algebraically closed. 
Let $\beta_1, \ldots, \beta_m \in X$ be such that $\cO_{X,\beta_i}$ is a valuation ring for $i =1, \ldots, m$. Let $Y \to X$ be the normalisation of a finite \'etale map $Y^* \to X^*$ that is tame at $\beta_1, \ldots, \beta_m$. Then there exist a $\lambda$ and an open $U \subseteq X_\lambda$ containing the images of $\beta_1, \ldots, \beta_m$ and such that, for every $\zeta \in \Spa(X_\lambda^*, X_\lambda)$ centred in $U$ (i.e. $\zeta \in \Spa(U^*,U)$), the cover $Y_\lambda \to X_\lambda$ is tame at $\zeta$. 
\end{lem}

\begin{proof}
Applying \cref{lem:tameness-conditions}~\ref{lemitem:tame-abhyankar} to the valuation rings $\cO_{X,\beta_i}$ and spreading out, we see that for each $i=1,\dots, m$ there exist
\begin{itemize}
    \item 
    an affine scheme $U_i$ with an \'etale map $U_i\to X$ whose image contains $\beta_i$,
    \item 
    elements $g_1,\dots, g_r\in \cO(U_i)$ whose images in $\cO(U_i^*)$ are invertible,
    \item 
    an integer $n\geq 1$ which is invertible on $X$;
    \item 
    a finite set $J_i$, and  
    \item 
    a surjective map of \'etale covers of $U_i^*$
    \[
        V_i = \coprod_{j \in J_i} 
        \Spec(\cO(U_i)[1/\pi][x_1,\dots, x_r]/(x_1^{n} - g_1, \dots, x_r^{n} - g_r)) \la Y^*\times_{X^*} U^*_i. 
    \]
\end{itemize}
Picking $\lambda$ large enough, we may assume that the $U_i\to X$, the functions $g_j\in \cO(U_i)$, and the map $V_i\to Y^*\times_{X^*}U_i^*$ are respectively the base changes to $X$ of \'etale maps $U_{i,\lambda}\to X_\lambda$, functions $g_j\in \cO(U_{i,\lambda})$ whose images in $\cO(U_{i,\lambda}^*)$ are invertible, and a surjection of \'etale covers of $U_{i,\lambda}^*$ from $V_{i,\lambda}$ (defined by the same formula over $U_{i,\lambda}$) onto the base change of $Y^*_\lambda$. It now suffices to take $U \subseteq X_\lambda$ to be the union of the images of the \'etale maps $U_{i,\lambda}\to X_\lambda$. Indeed, for every map $\Spec(V)\to U$ from a valuation ring $V$ which sends the generic point to $U^*$, the base change of $Y^*_\lambda \to X^*_\lambda$ to ${\rm Frac}(V)$ will be tamely ramified by the characterisation of  \cref{lem:tameness-conditions}~\ref{lemitem:tame-abhyankar}.
\end{proof}

\begin{lem} \label{lem:generic-points-ugggh}
    In the situation fixed after  \cref{lem:straightforward-approx}, let $\bar x\to X$ be a geometric point. Then for every $\lambda$, every open subset of the strict henselisation $X_{\lambda,(\bar x)}$ containing the images of the generic points of $X_{(\bar x)}\setminus (X_{(\bar x)})^*$ contains all generic points of $X_{\lambda,(\bar x)} \setminus (X_{\lambda,(\bar x)})^*$.
\end{lem}

\begin{proof}
Since $X\to W$ is \'etale, the induced map $X_{(\bar x)}\to W_{(\bar x)}$ is an isomorphism. We may therefore assume $X=W$ and $X_\lambda=W_\lambda$. Let us find the generic points of $W\setminus W^*$. Since $R_\lambda$ is a regular local ring in which $V(\pi)_{\rm red}$ is an snc divisor, there exist local parameters $x_1,\dots, x_d$ (where $d=\dim(R_\lambda)$), integers $a_1,\dots,a_r\geq 1$ with $1 \leq r\leq d$ and a unit $u \in R_{\lambda}^\times$ such that
\[
    \pi=u \cdot x_1^{a_1}\cdot\ldots\cdot x^{a_r}_r. 
\]
In this case, $W\setminus W^*=V(\pi)$ is the underlying topological space of the spectrum of 
\[
    R_\lambda[T_1,\dots, T_n]/(\pi - T_1\cdot\ldots\cdot T_s,\pi) = R_\lambda[T_1,\dots, T_n]/(x_1^{a_1}\cdot\ldots\cdot x^{a_r}_r, T_1\cdot\ldots\cdot T_s),
\]
and the minimal primes of this ring 
are of the form 
\[
\mathfrak{p}_{ij} = (x_i, T_j), \qquad i \leq r, \ j \leq s
\]
(if $s=0$, then this simplifies to $\mathfrak{p}_i=(x_i)$). The minimal primes of the ring 
\[
K^\circ[T_1,\dots, T_n]/(\pi, T_1\cdot\ldots\cdot T_s)
\]
are of the form 
\[
\mathfrak{q}_j = \sqrt{(\pi, T_j)},  \qquad j \leq s
\]
(if $s=0$ then this is just $\mathfrak{q} = \sqrt{(\pi)}$). Note that $x_i$ is contained in the maximal ideal of $K^\circ$ (since $R_\lambda \to K^\circ$ is local) that can be described as $\sqrt{(\pi)}$. We see that $\mathfrak{p}_{ij}\subseteq \mathfrak{q}_j$, resp.~$\mathfrak{p}_i \subseteq \mathfrak{q}$ if $s=0$. In other words, every generic point of $W_\lambda\setminus W^*_\lambda$ specialises to the image of a generic point of $W\setminus W^*$, and thus every open subset of $W_\lambda$ containing the images of the generic points of $W\setminus W^*$ contains all generic points of $W_\lambda\setminus W^*_\lambda$.

The computation shows that the irreducible components $V(\mathfrak{p}_{ij})$, resp.~$V(\mathfrak{p}_i)$ if $s=0$, of $W_\lambda\setminus W^*_\lambda$ are regular and in particular geometrically unibranch. It follows that the irreducible components of $W_{\lambda,(\bar x)} \setminus (W_{\lambda,(\bar x)})^*$ are in bijection with the irreducible components of $W_\lambda\setminus W^*_\lambda$ which pass through the image of $\bar x$ in $W_\lambda$, i.e.\ with those $\mathfrak{p}_{ij}$ such that $x_j(\bar x)=0=T_i(\bar x)$. Analogously, the irreducible components of $W_{(\bar x)}\setminus W^*_{(\bar x)}$ correspond to $\mathfrak{q}_i$ such that $T_i(\bar x) = 0$. Since $R_\lambda\to K^\circ$ is local and $\pi(\bar x)=0$, we have $x_j(\bar x)=0$ for all $j\leq r$. (The argument in the case $s=0$ is similar.) Again, it follows from this enumeration that every generic point $\zeta_{ij}$ of $W_{\lambda,(\bar x)} \setminus (W_{\lambda,(\bar x)})^*$ 
specialises to the image of a generic point $\beta_i$ of $W_{(\bar x)}\setminus W^*_{(\bar x)}$, and the result follows.
\end{proof}

We can now finish the proof of \cref{thm:scheme-Abhyankar} started earlier on page~\pageref{proof:scheme-Abhyankar}.

\begin{proof}[Proof of \cref{thm:scheme-Abhyankar} (part two)]
Assume that $K$ is algebraically closed and \ref{propitem:Abhy-tame-boundary} holds; we wish to show \ref{propitem:Abhy-ket} holds. Since the question is \'etale local on $X$, we pick a geometric point $\bar x\to X$ and show that it holds after base change to the strict henselisation  $X_{(\bar x)}$. We may assume that $X$ is affine and that there exists a pseudouniformiser $\pi$ and an \'etale map 
\[
    X\la W=\Spec(K^\circ[T_1,\dots, T_n]/(\pi-T_1\cdot\ldots\cdot T_s)).
\]
Apply \cref{lem:Zavyalov} and \cref{lem:straightforward-approx} to obtain the systems $Y^*_\lambda\to X_\lambda^*\to X_\lambda\to W_\lambda$ over $S_\lambda=\Spec(R_\lambda)$, as in the situation fixed after \cref{lem:straightforward-approx}. Apply \cref{lem:locus of tameness} with $\beta_i$ corresponding to the components of $V(\pi)$ to obtain a specific $\lambda$ and the open subset $U$, and then apply \cref{lem:generic-points-ugggh} to the preimage of $U$ in $X_{\lambda,(\bar x)}$. We deduce that the base change of $Y_\lambda\to X_\lambda$ to $(X_\lambda)_{(\bar x)}$ is tame at the generic points of $X_{\lambda,(\bar x)} \setminus (X_{\lambda,(\bar x)})^*$. Since $X_{\lambda,(\bar x)}$ is log regular with respect to the log structure induced by the open subset $(X_{\lambda,(\bar x)})^*$, we deduce from the classical purity \cref{thm:GR-log-Abhyankar} that the normalisation $Y_\lambda\to X_\lambda$ is Kummer \'etale at $\bar x\to X_\lambda$ (where we endow every scheme $Z\in \{S_\lambda, W_\lambda, X_\lambda, Y_\lambda \dots\}$ with the log structure induced by the open subset $Z^* = D(\pi)$). 

Let $Y'\to X$ be the map of log schemes obtained by base change of $Y_\lambda\to X_\lambda$ along the map of log schemes $X\to X_\lambda$. The map $Y'\to X$ is Kummer \'etale, and hence $Y'$ is smooth over $S$. It follows from \cite[5.3.1, 5.3.4, and 5.3.7]{AHLS-Log} that $Y'$ is $\eta$-normal, that $Y'\to X$ is finite, and that the log structure on $Y'$ is the one induced by the open subset $D(\pi)\subseteq Y'$. But then $Y'=Y$, and the result is proved.
\end{proof}

In order to deduce a version of purity for semistable formal schemes, we need an algebraisation-type result. 

\begin{lem}[{\cite[Corollaire, p.~578 + Remarque~2(c), p.~588]{Elkik}}] \label{Elkik-theorem} 
    Let $A$ be a ring and $\pi\in A$ a nonzerodivisor such that $(A, (\pi))$ is a henselian pair. Let $\widehat{A}$ be the $\pi$-adic completion of $A$. Then the morphism $A[1/\pi]\to \widehat{A}[1/\pi]$ induces an equivalence between finite \'etale $A[1/\pi]$-algebras and finite \'etale $\widehat{A}[1/\pi]$-algebras.
\end{lem}


\begin{cor} \label{finite-etale-Huber-pair}
    Let $(A,A^+)$ be a uniform Tate Huber pair, and let $\pi \in A^+$ be a topologically nilpotent unit.
    For every finite \'etale morphism $Y \to \Spa(A,A^+)$, we can find a homomorphism of Huber pairs
    $(A,A^+) \la (A_1,A_1^+)$
    such that 
    \begin{enumerate}[(i)]
    \item 
    $A^+ \to A_1^+$ is \'etale and $A_1 = A_1^+ \otimes_{A^+} A$, 
    \item 
    the induced map $A^+/\pi A^+\isomto A_1^+/\pi A_1^+$ is  an isomorphism, and 
    \item 
    there is a finite \'etale homomorphism of Huber pairs $(A_1,A_1^+) \la (B_1,B_1^+)$ and a cartesian square
    \[ 
    \begin{tikzcd}
        \Spa(B_1,B_1^+) \ar[d] \ar[r,"\sim"] & Y \ar[d] \\
        \Spa(A_1,A_1^+) \ar[r,"\sim"] & \Spa(A,A^+) . 
    \end{tikzcd}
    \]
    (Note that by (i) and (ii) the bottom map and therefore also the top map are isomorphisms.)
    \end{enumerate}
\end{cor}

\begin{proof}
    Combining \cite[1.4.4]{HuberBook} with \cite[Corollary~1.7.3(iii)]{HuberBook}, we know that finite étale morphisms to $\Spa(A,A^+)$ are in one to one correspondence with finite étale homomorphisms
    \[
     (\widehat{A},\widehat{A}^+) \longrightarrow (B,B^+)
    \]
    of Huber pairs.
    So $Y \to \Spa(A,A^+)$ comes from a Huber pair $(B_3,B_3^+)$ over $(\widehat{A},\widehat{A}^+)$ such that $\widehat{A} \to B_3$ is finite étale and~$B_3$ is the integral closure of~$\widehat{A}^+$ in~$B_3^+$.
    We consider the henselisation $(A^+_h,\pi)$ of the pair $(A^+,\pi)$ and set $A_h=A^+_h[1/\pi]$.
    We thus have homomorphisms of Huber pairs
    \[
     (A,A^+) \longrightarrow (A_h,A^+_h) \longrightarrow (\widehat{A},\widehat{A}^+) 
    \]
    and both of them induce isomorphisms on completions, hence on adic spectra.
    Moreover, this completion is just the $\pi$-adic completion of the ring of integral elements (here we use the uniformity assumption).
    By the generalised version of Elkik's theorem (\cref{Elkik-theorem}) the finite \'etale homomorphism
    \[
     \widehat{A} \longrightarrow B_3
    \]
    is the base change of a unique finite \'etale homomorphism
    \[
     A_h \longrightarrow B_2.
    \]
    Now we can write $A^+_h$ as a colimit of étale homomorphisms
    \[
     A^+ \longrightarrow A_1^+
    \]
    inducing an isomorphism modulo~$\pi$.
    Inverting~$\pi$ we obtain a presentation of~$A_h$ as a colimit of \'etale $A$-algebras~$A_1$.
    The finite \'etale map $A_h \to B_2$ comes via base change from some finite \'etale map
    \[
     A_1 \longrightarrow B_1.
    \]
    Taking~$B_1^+$ to be the integral closure of~$A_1^+$ in~$B_1$, we obtain the desired finite \'etale homomorphism
    \[
     (A_1,A_1^+) \longrightarrow (B_1,B_1^+). \qedhere
    \]
\end{proof}

In order to state the version of \cref{thm:scheme-Abhyankar} for semistable formal schemes, we first recall the construction of the normalisation of a formal scheme $\fX$ in a finite cover $Y$ of its generic fibre, see \cite[Lemma~4.1]{AchingerLaraYoucis2022}. Let $\fX$ be an $\eta$-normal (see \cref{defi:eta-normal}) tft formal scheme over $K^\circ$ with generic fibre $X = \fX_{\rm rig}$, and let $f\colon Y\to X$ be a finite morphism of rigid spaces with $Y$ reduced. Then there exists a unique (up to a unique isomorphism) $\eta$-normal formal model $\fY$ of $Y$ together with a finite map $\mathfrak{f}\colon \fY\to \fX$ with $\mathfrak{f}_\rig = f$. We call it the {\bf $\eta$-normalisation} of $\fX$ in $Y$. It can be characterised uniquely by the following two properties.
\begin{enumerate}[(i)]
    \item Its formation commutes with \'etale base change (see \cite[Appendix~A]{AchingerLaraYoucis2022}).
    \item If $\fX = \Spf(\mathcal{A})$ for a tft $K^\circ$-algebra $\mathcal{A}$, and $Y = \Spa(B)$ for a finite algebra $B$ over $A = \mathcal{A}_K$, then $\fY = \Spf(\mathcal{B})$ where $\mathcal{B}$ is the integral closure of $\mathcal{A}$ in $B$. 
\end{enumerate}
Since every semistable formal scheme is $\eta$-normal (see \cref{prop:log-sm-vertical}), the above construction applies in our situation of interest. We warn the reader that the construction of $\mathfrak{Y}$ (or more precisely the finiteness of $\mathcal{A}\to\mathcal{B}$ in (ii)) relies on our assumptions on $K$, see \cite[Remark~4.3]{AchingerLaraYoucis2022}.

The analogue of \cref{lem:local-ring-val-ring} for formal schemes is the fact (featured in \cref{rmk:sp-facts}) that for an $\eta$-normal tft formal scheme over $K^\circ$ and a generic point $\beta$ of $\fX_k$, the local ring $\cO_{\fX,\beta}$ is a rank one valuation ring. If $X$ is its rigid generic fibre, let us call a point $\zeta\in X$ a {\bf boundary point} of $\fX$ if its specialisation ${\rm sp}(\zeta)$ is a generic point of the special fibre $\fX_k$. Then the set of boundary points maps bijectively onto the set of generic points of $\fX_k$. Moreover, if $\beta = {\rm sp}(\zeta)$, then $\cO_{\fX,\beta}=\cO_{X,\zeta}^+=k(\zeta)^+$, and hence $k(\zeta)^\succ = k(\beta)$.

\begin{thm}[Purity, formal scheme version] \label{thm:Abhyankar-formal-more-precise}
  Let $\fX$ be a semistable formal scheme over $K^\circ$ and let $X=\fX_{\rm rig}$.  Assume that $K$ is either algebraically closed or discretely valued. 
 Let $Y\to X$ be a finite \'etale map and let $\fY\to \fX$ be the $\eta$-normalisation of $\fX$ in $Y$. The following are equivalent:
  \begin{enumerate}[(a)]
    \item \label{propitem:f-Abhy-tame} 
        $Y\to X$ is tame relative to $X$,
    \item \label{propitem:f-Abhy-tame-boundary}
        $Y\to X$ is tame above every boundary point of $\fX$, 
    \item \label{propitem:f-Abhy-ket}
        the morphism $\fY^{\rm log}\to\fX^{\rm log}$ is Kummer \'etale.
  \end{enumerate}
\end{thm}


\begin{proof}
The implication \ref{propitem:f-Abhy-tame}$\Rightarrow$\ref{propitem:f-Abhy-tame-boundary} is obvious and  \ref{propitem:f-Abhy-ket}$\Rightarrow$\ref{propitem:f-Abhy-tame} follows as in the scheme case (see the proof of \cref{thm:scheme-Abhyankar}). It remains to show that  \ref{propitem:f-Abhy-tame-boundary} implies \ref{propitem:f-Abhy-ket} by reducing to the case of schemes. Since the assertion is \'etale local on $\fX$, we may assume that $\fX$ is affine and \'etale over 
\[
    \mathfrak{W} = \Spf(\widehat{\mathcal{R}}), \qquad \mathcal{R} = K^\circ[T_1, \ldots, T_n]/(\pi - T_1\cdot\ldots\cdot T_s),
\]
where $\mathcal{R}$ is given the $\pi$-adic topology. By \cite[Corollary~1.7.3(iii)]{HuberBook} we find an \'etale ring homomorphism $\mathcal{R}\to \mathcal{A}$ such that $\fX = \Spf(\mathcal{A})$ (where $\mathcal{A}$ is endowed with the $\pi$-adic topology). Using \cref{finite-etale-Huber-pair} we can replace $\mathcal{A}$ with an \'etale $\mathcal{A}$-algebra with the same formal spectrum so that $Y\to X$ is induced by a finite \'etale morphism $\mathcal{A}[1/\pi]=A\to B$.

Consider the morphism of schemes of finite type over $K^\circ$
\[ 
    \Spec(B) \la \Spec(\mathcal{A}).
\]
Here $\Spec(\mathcal{A})$ is \'etale over $\Spec(\mathcal{R})$ and hence semistable, and $\Spec(B)\to \Spec(A) = \Spec(\mathcal{A}[1/\pi])$ is finite \'etale. Moreover, since $|\Spec(\mathcal{A}/\pi \mathcal{A})| = |\fX|$, the boundary points of $\fX$ are in bijection with the generic points of $\Spec(\mathcal{A}/\pi \mathcal{A})$. 

Let $\beta \in \Spec(\mathcal{A})$ be a generic point of $\Spec(\mathcal{A}/\pi \mathcal{A})$. Then $\cO_{\Spec(\mathcal{A}), \beta}$ is a valuation ring by \cref{lem:local-ring-val-ring}. We claim that $\Spec(B)\to \Spec(A)$ is tame above $\beta$ in the sense of \cref{thm:scheme-Abhyankar}. We have an extension of $\pi$-adically separated valuation rings 
\[ 
    \cO_{\Spec(\mathcal{A}), \eta} \la \cO_{\fX, \eta} = \cO_{X, \nu}^+ = k(\nu)^+.
\]
Here $\eta\in \fX$ denotes ``the same'' point on the formal scheme, and $\nu\in X = \fX_{\rm rig}$ is its unique preimage under the specialisation map. This map of valuation rings induces an isomorphism on $\pi$-adic completions, as for every $f\in \mathcal{A}$ the map $\mathcal{A}[f^{-1}]\to \widehat{\mathcal{A}}\langle f^{-1}\rangle = \widehat{\mathcal{A}[f^{-1}]}$  induces an isomorphism mod $\pi^n$ for all $n$. Passing to the colimit over all $f\in \mathcal{A}$ with $f(\beta)\neq 0$ we obtain the assertion. Therefore, tame extensions of $\cO_{\Spec(\mathcal{A}), \eta}$ correspond to tame extensions of $k(\nu)^+$.  

We are now in position to apply \cref{thm:scheme-Abhyankar} to $\Spec(B) \to \Spec(\mathcal{A})$, obtaining that if $\mathcal{B}$ is the integral closure of $\mathcal{A}$ in $B$, then $\Spec(\mathcal{B})\to\Spec(\mathcal{A})$ is Kummer \'etale with respect to the standard log structures. Let $\fY' = \Spf(\mathcal{B})$. Then $\fY'\to \fX$ is Kummer \'etale, and hence by \cref{prop:log-sm-vertical} we have that $\fY'$ is $\eta$-normal and it has the standard log structure. Clearly we have $\fY'_\rig = \Spa(B, \mathcal{B}) = Y$, and hence $\fY'\simeq \fY$. Thus $\fY\to \fX$ is Kummer \'etale, as desired.
\end{proof}

\begin{cor}
\label{cor:generic fibre and eta normalisation}
    Let $\fX$ be a semistable formal scheme over $K^\circ$ endowed with the standard log structure, and let $X=\fX_{\rig}$ be the rigid generic fibre.  Assume that $K$ is either algebraically closed or discretely valued. 
    The functor $\fY \mapsto Y = \fY_{\rm rig}$ induces an equivalence of Galois categories
    \[
    \FEt_\fX \longrightarrow \FEtt_{X/X} =\FEtt_X
    \]
    between finite Kummer \'etale covers of the log scheme $\fX$ and \'etale covers of $X$ which are tame (relative to $X$, see \cref{def:tame-rig}).
    The inverse functor is given by $\eta$-normalisation endowed with the standard log structure.
\end{cor}

\begin{proof}
Combining the four assertions of \cref{prop:log-sm-vertical},
if $\fY\to \fX$ is a Kummer \'etale map, then $\fY$ is $\eta$-normal and its log structure is the standard log structure. It now follows from \cref{thm:Abhyankar-formal-more-precise}~\labelcref{propitem:f-Abhy-ket}$\Rightarrow$\labelcref{propitem:f-Abhy-tame} that the generic fibre functor sends $\FEt_\fX$ into $\FEtt_{X}$.  Similarly, it follows from \cref{thm:Abhyankar-formal-more-precise}~\labelcref{propitem:f-Abhy-tame} $\Rightarrow$\labelcref{propitem:f-Abhy-ket} that the $\eta$-normalisation functor sends $\FEtt_{X}$ to $\FEt_\fX$. The two functors are mutual inverses. 
\end{proof}


%%%%%%%%%%%%%%%%%%%%%%%%%%%%%%%%%%%%%%%%%%%
\subsection{Comparison of tameness conditions}
\label{ss:tameness-comparison-pre}

The goal of this subsection is to prove an auxiliary result comparing tameness at infinity on the special and generic fibres of a formal scheme locally of finite type over~$K^\circ$. A finite \'etale map $f\colon \fY\to \fX$ of formal schemes locally of finite type over $K^\circ$ induces a finite \'etale map of $k$-schemes $f_k\colon \fY_k\to \fX_k$ and a finite \'etale map $f_{\rm rig}\colon \fY_{\rm rig}\to \fX_{\rm rig}$ of rigid spaces over $K$. The question is whether 
\[
    \text{$f_k$ is tame relative to $k$} \quad\overset{?}{\Leftrightarrow}\quad \text{$f_{\rm rig}$ is tame relative to $K$.}
\]
In \cref{prop:tameness-comparison-etale}, we confirm this under the assumption that the formal scheme $\fX$ is $\eta$-normal and that the scheme $\fX_k$ has the \emph{cdh uniformisation property} (\cref{def:cdh-unif}). This result is the key ingredient in the subsequent \S\ref{ss:tameness-comparison}, in which we prove a similar statement regarding finite Kummer \'etale covers of log formal schemes. 

\begin{rmk}
Let us explain the principal mechanism used to compare these two tameness conditions.
If we want to check tameness of $\fY_\rig \to \fX_\rig$ relative to~$K$, we need to consider a root triple $y = (y^\circ,k(y)^+,y^\succ)$ of~$\fY_\rig$ with image $x = (x^\circ,k(x)^+,x^\succ)$ in~$\fX_\rig$ and see whether the extension of valued fields $(k(y),k(y)^+)/(k(x),k(x)^+)$ is tamely ramified.
The valuation ring~$k(x)^+$ is the composition of $k(x^\circ)^+$ and a valuation ring~$V_x$ on the specialisation field~$k(x^\circ)^\succ$ and similarly for~$k(y)^+$ (call the corresponding valuation ring~$V_y$).
Since $\fY_\rig \to \fX_\rig$ is the generic fibre of the \'etale map $\fY\to \fX$, we already know that $(k(y^\circ),k(y^\circ)^+)/(k(x^\circ),k(x^\circ)^+)$ is unramified (and in particular tame).
But tameness behaves well with respect to composite valuations (\cref{composition-tame}), so $(k(y),k(y)^+)/(k(x),k(x)^+)$ is tame if and only if $(k(y^\circ)^\succ,V_y)/(k(x^\circ)^\succ,V_x)$ is tame.

The valued field $(k(x^\circ)^\succ,V_x)$ defines a point~$x_1$ on $\Spa(\fX_k,k)$ by restricting~$V_x$ to the residue field $k(\spe x^\circ)$ of the image of~$x^\circ$ under the specialisation map $\spe \colon X \to \fX_k$.
In a similar way we define the point $y_1  \in \Spa(\fY_k,k)$. 

If $\fY_k \to \fX_k$ is tame relative to~$k$, the extension of valued fields $(k(y_1),k(y_1)^+)/(k(x_1),k(x_1)^+)$ is tamely ramified.
The tricky question is whether this implies that $(k(y^\circ)^\succ,V_y)/(k(x^\circ)^\succ,V_x)$ is tame.
For the converse direction we have to answer the even more difficult question: does the tameness of $(k(y^\circ)^\succ,V_y)/(k(x^\circ)^\succ,V_x)$ imply the tameness of $(k(y_1),k(y_1)^+)/(k(x_1),k(x_1)^+)$?
As a direct conclusion this is not true in general. 
There might be a nontrivial residue extension $k(x^\circ)^\succ/k(x_1)$ and the wild ramification of $k(y_1)/k(x_1)$ might be killed when base changing the extension to $k(x^\circ)^\succ$.
However, we will use that under an extra assumption (the cdh uniformisation property) it is enough to test tameness only at points of a special type (\cref{tameness-support-generic-point}).
For these points we can ensure that $k(x^\circ)^\succ = k(x_1)$.

We start this subsection with the construction of these special points.
\end{rmk}

In the two statements below (\cref{lem:valuation-on-regular} and \cref{lem:valuation-on-cdh-unif}), we use the following terminology. Let $\xi\leadsto x$ be a specialisation of points on a scheme $X$ and let $k(\xi)^+\subseteq k(\xi)$ be a valuation subring. We say that $k(\xi)^+$ has {\bf residue field $k(x)$} if the map $\cO_{X,x}\to \cO_{X,\xi}\to k(\xi)$ has image in $k(\xi)^+$, the map $\cO_{X,x}\to k(\xi)^+$ is local, and induces an isomorphism of residue fields $k(x)\simeq k(\xi)^\succ$. Equivalently, there exists a (necessarily unique) map $\Spec(k(\xi)^+)\to X$ extending $\Spec(k(\xi))\to X$, mapping the closed point $s$ to $x$, and inducing an isomorphism $k(x)\simeq k(\xi)^\succ =k(s)$. 

\begin{lem} \label{lem:valuation-on-regular}
    Let $\xi\leadsto x$ be a specialisation on a scheme $X$ and let $Z=\overline{\{\xi}\}$ with the reduced subscheme structure. Suppose that the local ring $\cO_{Z,x}$ is regular. Then there exists a valuation subring $k(\xi)^+\subseteq k(\xi)$ with residue field $k(x)$.
\end{lem}

\begin{proof}
We may replace $X$ with $\Spec(\cO_{Z,x})$ and thus reduce to the case where $X=Z$ is the spectrum of a regular local ring $R=\cO_{X,x}$ with generic point $\xi$ and closed point $x$. Let $f_1, \ldots, f_d$ be a regular sequence of parameters in $R$, which yields a sequence of specialisations
\[
 \xi = x_0 \rightsquigarrow x_1 \rightsquigarrow \cdots \rightsquigarrow x_d = x,
\]
where $x_i$ corresponds to the prime ideal $(f_1, \dots, f_i)$ and each $\cO_{\overline{\{x_{i-1}\}},x_i}$ is a discrete valuation ring with uniformiser $f_i$, fraction field $k(x_{i-1})$, and residue field $k(x_i)$.
Let $k(\xi)^+$ be the valuation ring of $k(\xi)$ obtained as the composition of all these discrete rank one valuation rings.
It is discrete of rank~$d$ and its residue field $k(\xi)^\succ$ equals $k(x_d) = k(x)$.
\end{proof}

We can generalise the above result to certain non-regular schemes. For this we need the following definition.

\begin{defi} \label{def:cdh-unif}
    A quasi-excellent scheme $X$ has the {\bf cdh uniformisation property} if for every $x\in X$ and every irreducible component $Z\subseteq X$ containing $x$ (with the reduced scheme structure) there exists a dominant finite type morphism $f\colon Z'\to Z$ with $Z'$ integral and regular, and a point $x'\in f^{-1}(x)\subseteq Z'$ with $k(x)=k(x')$.
\end{defi}

It is easy to see that the cdh uniformisation property can be checked Zariski locally.
Moreover, if $Y\to X$ is a smooth map and $X$ has the cdh uniformisation property, then so does $Y$.
However, the property is not étale local as the following example shows.

\begin{ex} \label{ex:nodal-curve-noncdh}
    Consider the nonsplit nodal cubic over $\mathbb{R}$:
    \[
        X = \Spec(\mathbb{R}[x,y]/(y^2 + x^2 - x^3))
    \]
    and let $P=(x=y=0)\in X$ be the node. The normalisation $X'\to X$ is given by $X'=\Spec(\mathbb{R}[t])$ where $x=t^2+1$ and $y=t^3+t=tx$, and the point $P'=(t^2+1=0)\in X'$ is the unique preimage of $P$. Then $X$ does not have the cdh uniformisation property. Indeed, suppose that $f\colon Y\to X$ is a dominant finite type map with $Y$ regular and $Q\in Y$ satisfies $f(Q)=P$ and $k(Q)=k(P)=\mathbb{R}$. Since $Y$ is normal, $f$ factors through $X'$, and then $Q$ maps to $P'$, so that $k(Q)\supseteq k(P')=\mathbb{C}$, contradiction. However, the base change $X_\mathbb{C}$ does have the cdh uniformisation property, as exhibited by its normalisation $X'_\mathbb{C}\to X_\mathbb{C}$.
\end{ex}

The main point about schemes with the cdh uniformisation property is that we can dominate their local rings with valuation rings with the same residue field.

\begin{lem} \label{lem:valuation-on-cdh-unif}
    Let $\xi\leadsto x$ be a specialisation on a scheme $X$. Suppose that the scheme $Z = \overline{\{\xi\}}$ (with the reduced subscheme structure) has the cdh uniformisation property. Then there exists a valuation subring $k(\xi)^+\subseteq k(\xi)$ with residue field $k(x)$. 
\end{lem}

\begin{proof}
Let $Z'\to Z$ and $x'\in Z'$ be as in \cref{def:cdh-unif}. Let $\xi'\in Z'$ be the generic point, which maps to $\xi$. By \cref{lem:valuation-on-regular} there exists a valuation subring $k(\xi')^+\subseteq k(\xi')$ with residue field $k(x')=k(x)$. Then $k(\xi')^+\cap k(\xi)\subseteq k(\xi)$ is a valuation subring with the same residue field. 
\end{proof}

For our purposes, the cdh uniformisation property is important since it ensures that tameness can be checked at valuations with generic support.

\begin{lem}  \label{tameness-support-generic-point}
    Let $f\colon Y\to X$ be a finite \'etale  morphism of quasi-excellent schemes over a base $S$. Suppose that $X$ has the cdh uniformisation property. If 
    \[
        \Spa(f) \colon \Spa(Y,S) \la \Spa(X,S)
    \]
    is tame at all points of $\Spa(Y,S)$ whose support is a generic point of~$Y$, then $\Spa(f)$ is tame everywhere.
\end{lem}

\begin{proof}
Suppose that $\Spa(f)$ is tame at all points with generic support.
We want to show that $\Spa(f)$ is tame at $(y,k(y)^+)\in \Spa(Y, S)$ (we omit the centre of $k(y)^+$ on $S$ from the notation). Let $(x, k(x)^+)$ be its image in $\Spa(X, S)$ and let $x'$ be a generic point of $X$ specialising to $x$. By \cref{lem:valuation-on-cdh-unif}, there exists a valuation subring $k(x')^+\subseteq k(x')$ with residue field $k(x)$. We thus have a map $V=\Spec(k(x')^+)\to X$ sending the generic point to $x'$, the special point to $x$, both without changing the residue field. In particular, the fibre $Y_x$ containing $y$ can be regarded as a closed subscheme of $W = Y\times_X V$. The ring $k(y')^+ = \cO_{W,y}$ is a valuation ring with residue field $k(y)$ and fraction field $k(y')$ for a unique $y'\in Y_{x'}$. Since $y\mapsto x$, the map $k(x')^+\to k(y')^+$ thus constructed is local, so that $(k(y'), k(y')^+)/(k(x'), k(x')^+)$ is an extension of valuation rings. Since $W\to V$ is finite \'etale, this extension is moreover unramified. Composing each of these valuation rings with the valuation rings in the extension 
\[
    (k(y')^\succ=k(y), k(y)^+)  /  (k(x')^\succ = k(x), k(x)^+)
\]
we obtain a point $(y', k(y')^+\times_{k(y)} k(y)^+)\in \Spa(Y, S)$ lying over $(x', k(x')^+\times_{k(x)} k(x)^+)\in \Spa(X, S)$. Since $Y\to X$ is flat, $y'$ is a generic point of $Y$, and thus by our assumption, the corresponding extension of valuation rings is tame. We are now in position to apply \cref{composition-tame}, which implies that $(k(y), k(y)^+)/(k(x),k(x)^+)$ is tame, so that $\Spa(f)$ is tame at $(y, k(y)^+)$.
\end{proof}


In accordance with \cref{def:tame-rig}, we call a finite \'etale  map of separated finite type formal schemes $\fY\to \fX$ {\bf tame relative to $K^\circ$} if the induced finite \'etale map of universal compactifications
\[
    \Spa(\fY^{\rm ad}, K^\circ) \la \Spa(\fX^{\rm ad}, K^\circ)
\]
is tame.
Note that~$\fX^{\rm ad}$ is indeed an adic space and satisfies the hypotheses of \cref{thm:Solodov-comp} (see \cref{rmk:hypotheses-compactification}), so the universal vertical compactification exists and similarly for~$\fY$.
The above definition is clearly local on~$\fX$.
If~$\fX$ is not necessarily separated and quasi-compact, we can define tameness of a finite étale morphism $\fY \to \fX$ by requiring it to be tame when restricted to all affine opens (for which the universal vertical compactification is known to exist).


\begin{prop} \label{prop:tameness-comparison-etale}
    Let $\fX$ be a formal scheme locally of finite type over $K^\circ$ and let $f \colon \fY\to \fX$ be a finite \'etale map.
    Consider the following conditions:
    \begin{enumerate}[(a)]
        \item \label{propitem:tameness-comparison-etale-a} The finite \'etale map of $k$-schemes $f_k \colon \fY_k\to\fX_k$ is tame relative to $k$;
        \item \label{propitem:tameness-comparison-etale-b}The map of formal schemes $f \colon \fY\to\fX$ is tame relative to $K^\circ$;
        \item \label{propitem:tameness-comparison-etale-c} The finite \'etale map of rigid spaces $f_\rig \colon \fY_{\rm rig}\to\fX_{\rm rig}$ is tame relative to $K$.
    \end{enumerate}
    Then \labelcref{propitem:tameness-comparison-etale-a} $\Leftrightarrow$ \labelcref{propitem:tameness-comparison-etale-b} $\Rightarrow$ \labelcref{propitem:tameness-comparison-etale-c}, and if $\fX$ is $\eta$-normal and $\fX_k$ has the cdh uniformisation property, then all three conditions are equivalent.
\end{prop}

\begin{proof}
The assertion is local on~$\fX$, so we may assume that~$\fX$ is separated and quasi-compact.
By \cref{cor:base-change-comp} the generic fibre of the induced map of universal compactifications $\Spa(\fY^{\rm ad}, K^\circ) \to \Spa(\fX^{\rm ad}, K^\circ)$ identifies with $\Spa(\fY_\rig, K) \to \Spa(\fX_\rig, K)$ and its special fibre is $\Spa(\fY_k,k) \to \Spa(\fX_k,k)$.
Since tameness is stable under base change, it is now clear that~\labelcref{propitem:tameness-comparison-etale-b} implies~\labelcref{propitem:tameness-comparison-etale-a} and~\labelcref{propitem:tameness-comparison-etale-c}.

Let us now show that~\labelcref{propitem:tameness-comparison-etale-a} implies~\labelcref{propitem:tameness-comparison-etale-b}.
The morphism $\Spa(\fY^{\rm ad},K^\circ) \to \Spa(\fX^{\rm ad},K^\circ)$ is finite étale by \cref{prop:tame-S-compactification}.
By assumption, it is tame at all points of the special fibre.
Every other point of $\Spa(\fY^{\rm ad},K^\circ)$ has a specialisation lying in the special fibre.
By the openness of the tame locus \cite[Corollary~4.4]{Hubner2021:AdicTameSite} it follows that $\Spa(\fY^{\rm ad}, K^\circ) \to \Spa(\fX^{\rm ad}, K^\circ)$ is tame.

Finally, let us deduce~\labelcref{propitem:tameness-comparison-etale-a} from~\labelcref{propitem:tameness-comparison-etale-c} in case~$\fX$ is $\eta$-normal and~$\fX_k$ has the cdh uniformisation property.
Suppose that $\Spa(\fY_\rig, K) \to \Spa(\fX_\rig, K)$ is tame.
Being étale over~$\fX_k$, the scheme $\fY_k$ inherits the cdh uniformisation property from~$\fX_k$. 
Therefore, by \cref{tameness-support-generic-point}, to show that $\Spa(\fY_k, k) \to \Spa(\fX_k, k)$ is tame, it suffices to test tameness at valuations whose support is a generic point of an irreducible component.
So let us take a valuation ring~$V_y$ on $k(\eta)$ for the generic point~$\eta$ of an irreducible component of $\fY_k$.
Since $\fX$ is $\eta$-normal, so is~$\fY$, and by \cref{rmk:sp-facts}(3) there is a (unique) point $y \in \fY_\rig$ with $\spe(y) = \eta$ for which
\[
k(y)^\succ = k(\eta).
\]
Let $\bar{y}$ be the composition of~$y$ with~$V_y$.
This is a point of~$\Spa(\fY_\rig,K)$.
We now map the whole constellation to~$\Spa(\fX_\rig,K)$:
\[
x = f(y),    \quad   \bar{x} = \Spa(f)(\bar{y}),   \quad   \eta_X = f(\eta),   \quad V_x = V_y \cap k(\eta_X).
\]
Then $k(x)^\succ = k(\eta_X)$ and $(k(\bar{x}),k(\bar{x})^+) = (k(x),k(x)^+) \succ (k(\eta_X),V_x)$.
We know by assumption that the extension of valued fields
\[
(k(\bar{y}),k(\bar{y})^+)/(k(\bar{x}),k(\bar{x})^+)
\]
is tame.
Therefore, \cref{composition-tame} implies that $(k(\eta),V_y)/(k(\eta_X),V_x)$
is tame, i.e.~$\Spa(f_k)$ is tame at~$(\eta,V_x)$.
\end{proof}


%%%%%%%%%%%%%%%%%%%%%%%%%%%%%%%%%%%%%%%%%%%
\subsection{Tame covers and Kummer covers (II): tameness at infinity}
\label{ss:tameness-comparison}

We consider a  strictly semistable formal scheme $\fX$ over~$K^\circ$ with rigid generic fibre~$\fX_\rig$ and special fibre~$\fX_k$.
When  $K$ is either algebraically closed or discretely valued, by \cref{cor:generic fibre and eta normalisation} 
we have equivalences
\[
 \FEt^t_{\fX_\rig} \simeq \FEt_\fX \simeq \FEt_{\fX_k},
\]
where $\FEt_\fX$ and $\FEt_{\fX_k}$ are the respective categories of finite Kummer \'etale covers.
In this subsection, we add tameness conditions at the boundary to the picture.
Let $\fY \to \fX$ be a finite (Kummer) étale morphism of log formal schemes, let $\fY_\rig \to \fX_\rig$ be its rigid generic fibre (a tame finite étale map of rigid spaces) and let $\fY_k \to \fX_k$ be its special fibre (a finite \'etale map of log schemes over the residue field~$k$). In \cref{prop:tameness-comparison} below, we shall prove that $\fY_\rig \to \fX_\rig$ is tame relative to $K$ (in the sense of \cref{def:tame-rig}) if and only if $\fY_k\to \fX_k$ is tame relative to $k$ (in the sense of \cref{def:tame-log}).
As a consequence we obtain equivalences of categories
\[
 \FEt^t_{\fX_\rig/K} \simeq \FEt_{\fX/K^\circ} \simeq \FEt_{\fX_k/k}.
\]
Ultimately, we shall deduce this result from \cref{prop:tameness-comparison-etale}. In order to apply it, we first need to be able to ensure that $\fX_k$ has the cdh uniformisation property.
We start our investigation with two preparatory lemmas.

\begin{lem} \label{lem:extend-valuative}
    Let $\alpha \colon P \to k$ be a morphism from an fs monoid to the multiplicative monoid of a field.
    Then there is a valuative monoid $V$ with $P \subseteq V \subseteq P^\gp$ such that~$f$ extends to a morphism $f_V \colon V \to k$.
\end{lem}

\begin{proof}
    Let~$F$ be the face $f^{-1}(k^\times) = P \setminus f^{-1}(0)$ of~$P$.
    By the universal property of localisation, we can extend~$f$ to the localisation $P-F$.
    Therefore, we may assume in the following that $f$ is local, i.e.\ that it maps all non-units to~$0$ (and all units to $k^\times$).

    By \cite[Proposition~I.2.4.1]{Ogus} there is a valuative monoid~$V$ dominating~$P$, i.e. $P \subseteq V \subseteq P^\gp$ such that $P \to V$ is local.
    We claim that~$f$ extends to any such monoid (in fact we do not need it to be valuative, it works with any monoid dominating~$P$).
    Restricting~$f$ to the units of~$P$ we obtain a group homomorphism
    \[
        f^\times \colon P^\times \longrightarrow k^\times
    \]
    and we want to extend it to~$V^\times$.
    Note that~$V^\times$ and~$P^\times$ are finitely generated abelian groups as they are both subgroups of the finitely generated abelian group~$P^\gp$.
    Moreover, the quotient $V^\times/P^\times$ is torsion-free by saturatedness of~$P$.
    We conclude that $V^\times/P^\times$ is free and $V^\times \simeq P^\times \oplus (V^\times/P^\times)$.
    Hence, we can easily extend~$f^\times$ to $f_V^\times \colon V^\times \to k^\times$.

    We now define a homomorphism of monoids $f_V \colon V \to k$  by setting $f_V|_{V^\times} = f_V^\times$ and sending all non-units to zero.
    This is indeed an extension of~$f$ because $P \to V$ is local and~$f$ maps all non-units to zero.
\end{proof}

\begin{lem} \label{lem:log-blowup-fibres}
    Let $X$ be an fs log scheme admitting a chart Zariski locally and let $f\colon X'\to X$ be a log blow-up.
    For every point $x\in X$ there exists an $x'\in f^{-1}(x)$ with $k(x')=k(x)$.
\end{lem}

\begin{proof}
Since the formation of log blowups is compatible with strict base change, we easily reduce to the case $X = \Spec(P\xrightarrow{\alpha} k)$ for a field $k$ and an fs monoid $P$. In this case, $X'$ is the blow-up corresponding to a finitely generated ideal subsheaf  $\overline{\mathcal{J}}\subseteq \overline{\cM}_X$. Since the latter is the constant sheaf with value $P/P^\times$ over $\Spec(k)$, the sheaf $\overline{\mathcal{J}}$ is also constant, and hence $X'\to X$ is the log blowup corresponding to an ideal $J\subseteq P$ (the preimage of $\overline{\mathcal{J}}(X)\subseteq P/P^\times$ in $P$). Using \cref{lem:extend-valuative} we can extend $\alpha \colon P \to k$ to a monoid homomorphism $V \to k$ for a valuative monoid~$V$ with $P \subseteq V \subseteq P^\gp$.
In~$V$ the ideal~$J$ becomes principal as all finitely generated ideals of valuative monoids are principal.
By the universal property of a log blow-up (see \cite[Definition~III.2.6.2]{Ogus}, note that this is applicable even though~$V$ is not fine by passing to a suitable fine submonoid of~$V$), the map
\[ 
    \Spec(V\to k) \longrightarrow \Spec(P\xrightarrow{\alpha}k) = X
\]
factors through the blow-up $X'$. This shows that $X'(k)\neq \emptyset$, as desired.
\end{proof}

\begin{prop} \label{lem:log-sm-cdh-unif}
    Let $S = \Spec(M\xrightarrow{\alpha} k)$ where $k$ is a field and $M$ a saturated monoid that is finitely generated over a sharp divisible valuative monoid~$V$, and where $\alpha^{-1}(k^\times)=\{0\}$.
    Let $f\colon X\to S$ be a log smooth morphism which admits a chart Zariski locally; that is, Zariski locally on $X$ there exists a smooth map of saturated monoids $M\to P$ and a strict \'etale morphism
    \[
        X \la \Spec(P\to k[P] \otimes_{k[M]} k).
    \]
    Then the underlying scheme $\underline{X}$ has the cdh uniformisation property.
\end{prop}

\begin{proof}
We are in the situation of \cref{setup:surgery}.
The normalisation map $\nu \colon X^{(0)} \to X$ induces isomorphisms on residue fields by \cref{lem:surgery-properties}.
Moreover, the log scheme~$X_\star^{(0)}$ (which has the same underlying scheme as~$X^{(0)}$ but with the compactifying log structure defined by the locally constant locus in~$X^{(0)}$) is fs and smooth over~$k$ (with the trivial log structure).
By toric resolution of singularities (see \cite[(10.4)]{KatoToricSingularities}) there exists a log blowup $Z \to X_\star^{(0)}$ such that~$\underline{Z}$ is regular.
Now $\underline{Z} \to \underline{X}_\star^{(0)} \to \underline{X}$ is surjective and every point $x \in \underline{X}$ can be lifted to a point $x' \in \underline{X}_\star^{(0)}$ with trivial residue field extension.
Moreover, from \cref{lem:log-blowup-fibres} we can further lift~$x'$ to $z\in \underline{Z}$  with $k(z)=k(x')$.
This shows the cdh uniformisation property.
\end{proof}

We now come back to the setting we are interested in:
we want to study tameness at the boundary of finite Kummer étale morphisms $\fY \to \fX$ of log formal schemes over~$K^\circ$.

\begin{prop} \label{prop:tameness-comparison}
    Let $\fX$ be a strictly semistable log formal scheme over $K^\circ$ and let $\fY\to\fX$ be a finite Kummer \'etale map.  Assume that $K$ is either algebraically closed or discretely valued. 
 Then the following are equivalent:
    \begin{enumerate}[(a)]
        \item the finite Kummer \'etale map of special fibres $\fY_k\to \fX_k$ is tame relative to $k$,
    
        \item the finite \'etale map of generic fibres $\fY_\rig\to\fX_\rig$ is tame relative to $K$.
    \end{enumerate}
\end{prop}

\begin{proof}
We set $M = \Gamma_K^+$. 
Since the assertion is Zariski local on $\fX$, we may assume by \cref{lem:strictly-semistable-formal-scheme-reduces-to-setup364} that there is a chart
\[
    \fX \la \Spf(P \to K^\circ \otimes_{k[M]}k[P])
\]
with $M \to P$ a smooth morphism. 
For a Kummer homomorphism of monoids $P \to Q$ we consider the cartesian squares
\begin{equation} \label{diagram:cover-P-Q}
 \begin{tikzcd}
     \fY_Q    \ar[r]  \ar[d]  
     & \fX_Q   \ar[r]  \ar[d]  &
     \Spf(Q \to K^\circ \otimes_{k[M]}k[Q])  \ar[d]  \\
     \fY      \ar[r]          
     & \fX     \ar[r]       &    \Spf(P \to K^\circ \otimes_{k[M]}k[P]).
 \end{tikzcd}
\end{equation}
Passing to special fibres we obtain an analogous diagram of log schemes over~$k$.
By \cref{cor:fKet-local-structure} we can choose $P \to Q$ in such a way that $\fY_{Q,k} \to \fX_{Q,k}$ is strict étale.
Since strict étaleness can be checked on the special fibre, also $\fY_Q \to \fX_Q$ is strict étale.
The morphism $(\fX_Q)_k \to \fX_k$ is automatically tame relative to~$k$ as it is a torsor under a finite abelian group scheme of degree prime to the residue characteristic. 
So in the base change of the diagram~(\ref{diagram:cover-P-Q}) to $k$, the left and middle vertical arrows are tame relative to~$k$ and therefore the upper left horizontal arrow $(\fY_Q)_k \to (\fX_Q)_k$ is tame relative to~$k$ if and only if the lower left horizontal arrow $\fY_k \to \fX_k$ is tame relative to~$k$ (see \cref{lem:tame-log-properties}\ref{lemitem:tame-log-properties-XQ}). The same reasoning applies to the generic fibres.

This reduces us to proving the proposition for $\fY_Q \to \fX_Q$ instead of $\fY \to \fX$.
In particular, we may assume that our morphism is strict étale.
Note that by construction $\fX_Q\to \Spf(K^\circ)$ admits a global chart such that the special fibre $\fX_{Q,k}^{\log}$ satisfies the \cref{setup:surgery}, in particular because the composition $M \to P \to Q$ is smooth. Therefore \cref{lem:log-sm-cdh-unif} applies, and hence its special fibre $\fX_{Q,k}$ has the cdh uniformisation property. 
By \cref{lem:tame-log-properties}\ref{lemitem:tame-log-properties-strict} we can thus forget about the log structure on~$\fX$ and just study finite étale morphisms $\fY \to \fX$ of $\eta$-normal formal schemes such that~$\fX_k$ has the cdh uniformisation property.
Then the result follows from \cref{prop:tameness-comparison-etale}.
\end{proof}

Combining with \cref{cor:generic fibre and eta normalisation}, we obtain the main result of this section:

\begin{cor} \label{cor:Abhy-infty-equiv}
    Let $\fX$ be a strictly semistable log formal scheme over $K^\circ$. Assume that $K$ is either algebraically closed or discretely valued. 
    The functor ``generic fibre'' 
    $\FEt_{\fX_k}\simeq \FEt_\fX\to \FEt_{\fX_\rig}$ induces equivalences
    \[
        \FEtt_{\fX_k/k} \simeq \FEtt_{\fX/K^\circ} \simeq \FEtt_{\fX_\rig/K}.
    \]
\end{cor}

\cref{lem:log-sm-cdh-unif} is a crucial ingredient in the proof of \cref{cor:Abhy-infty-equiv}. In \cref{lem:log-sm-cdh-unif} it is important that the chart exists Zariski locally and not only étale locally because the cdh uniformisation property is not étale local.
But this property was key to constructing sufficiently nice points where to test tameness.
This is the reason why we work with \emph{strictly} semistable formal models in this subsection. Note that the nodal cubic in \cref{ex:nodal-curve-noncdh} is the special fibre of a semistable but not strictly semistable formal scheme over $\mathbb{R}(\!(\pi)\!)$.

%%%%%%%%%%%%%%%%%%%%%%%%%%%%%%%%%%%%%%%%%%%
\subsection{Finite generation, finite presentation, and examples}
\label{ss:tamepi1-rig-fgfp}

We are now ready to prove the main results of this section regarding finite generation and finite presentation of $\pitame(X/K)$. 

\begin{thm} \label{thm:fg-pi1-rig}
    Let $K$ be a non-archimedean field whose tame Galois group $\pitame(K)$ is finitely generated. Let $X$ be a qcqs connected rigid analytic space over $K$ with geometric point $\bar x\to X$. 
    Then $\pitame(X/K,\bar x)$ is finitely generated.
\end{thm}

\begin{proof} 
Let $C$ be a completed algebraic closure of $K$. There exists a finite extension $K'\subseteq C$ of $K$ such that the connected components of $X_{K'}$ are geometrically connected (see \cite[Corollary 3.2.3]{Conrad:IrreducibleComponents}). By descent (\cref{cor:descent-of-fg-fett}), we may replace $X$ with one of those components and $K$ with $K'$ and hence assume that $X$ is geometrically connected. By \cref{cor:fg-alg-closed}, we may then replace $K$ with $C$ and $X$ with $X_C$. We have thus reduced to the case where $K$ is algebraically closed. 
   
Let $Y\to X$ be a surjection from a regular (and hence smooth, since $K$ is algebraically closed) qcqs rigid space $Y$, which exists by \cref{prop:h-locally-regular}. By Temkin's uniformisation result recalled as \cref{thm:Temkin-unif}, there exists an \'etale surjection $Z\to Y$ such that $Z$ is qcqs and admits a strictly semistable formal model~$\fZ$.

Let $\fW^{\rm log}$ be a connected component of $\fZ$ endowed with the standard log structure and $W\subseteq Z$ the corresponding connected component of $Z$. By \cref{cor:Abhy-infty-equiv} we have an isomorphism of profinite groups
\[ 
    \pitame(W/K) \simeq \pitame(\fW^{\rm log}_k/k).
\]
Combining this with \cref{thm:pi1tame log is fg}, we conclude that the tame fundamental group of every connected component of~$Z$ is finitely generated. Applying descent (\cref{cor:descent-of-fg-fett}) to the surjective map $Z\to X$, we deduce the finite generation of $\pitame(X/K)$.
\end{proof}

\begin{cor}
    In the setting of \cref{thm:fg-pi1-rig}, the prime-to-$p$ quotient of the étale fundamental group $\pi_1(X)$ is finitely generated (where~$p$ is the residue characteristic exponent of~$K$).
\end{cor}

\begin{proof}
The category of all étale covers of~$X$ whose connected components have Galois closures of degree prime to~$p$ is a full subcategory of $\FEtt_{X/K}$ satisfying the conditions of \cref{lem:Galois-subcategory-criterion}.
The corresponding fundamental group is thus a quotient of $\pitame(X/K)$ and inherits the property of being finitely generated.
\end{proof}

Let us take a look at Berkovich's tame fundamental group~$\pi_1^B(X/K)$ studied in \cite[\S6.3]{Berkovich:EtaleCohomology}.
In \cref{rmk:Berkovich-covers} we noted that $\pi_1^B(X/K)$ is a quotient of $\pitame(X/K)$.
Consequently, we can deduce from \cref{thm:fg-pi1-rig} that it is finitely generated:

\begin{cor}
    In the setting of \cref{thm:fg-pi1-rig}, Berkovich's tame fundamental group $\pi_1^B(X/K)$ is finitely generated.
\end{cor}

The method of proof of \cref{thm:fg-pi1-rig} is inexplicit, in that it does not provide a formula for $\pitame(X/K)$. Below, we show how $\pitame(X/K)$ can be computed in practice.  We use \cref{cor:Abhy-infty-equiv} and the results of \S\ref{ss:surgery} to give a partial result regarding the finite presentation of $\pitame(X/K)$ of a rigid space $X$ over an algebraically closed $K$. Afterwards, we use the same approach to compute some examples.

\begin{thm} \label{thm:fp-sometimes}
    Let $K$ be an algebraically closed non-archimedean field and let $X/K$ be the rigid generic fibre of a connected quasi-compact strictly semistable formal scheme $\fX/K^\circ$. Let $U\subseteq \fX_k$ be the smooth locus of the special fibre (with no log structure). Suppose that $U$ admits a projective snc compactification. Then $\pitame(X/K)$ is finitely presented. 
\end{thm}

\begin{proof}
By \cref{cor:Abhy-infty-equiv} we have $\pitame(X/K)\simeq \pitame(\fX_k/k)$.
The special fibre $\fX_k$ endowed with the restriction of the standard log structure of $\fX$ satisfies the assumptions in \cref{setup:surgery} by \cref{lem:strictly-semistable-formal-scheme-reduces-to-setup364}. We conclude that $\pitame(\fX_k/k)$ is finitely presented thanks to \cref{cor:surgery-fp}.
\end{proof}

\begin{rmk}
    Our methods show that in order to check tameness, it suffices to check a finite number of rank one and two points of the universal compactification. More precisely, in the situation of \cref{thm:fp-sometimes}, let $U\hookrightarrow \overline{U}$ be an snc compactification, let $\beta_1, \dots, \beta_n$ be the generic points of $\fX_k$, and let $\gamma_1, \dots, \gamma_r$ be the generic points of $\overline{U}\setminus U$. As noted in \cref{rmk:sp-facts}, for each $i=1,\dots, n$ there is a rank one point $\zeta_i$ of $X$, the unique point specialising to $\beta_i$. Moreover, for each $i=1,\dots, r$, there is a rank two point $\xi_i$ of the universal compactification $\overline{X}$, whose valuation is described as the composition of the valuation corresponding to the irreducible component $Z$ of $\fX_k$ such that $\gamma_i$ lies in the closure of $Z \cap U$ with the discrete valuation on $k(Z) = k(Z \cap U)$ with valuation ring $\cO_{\overline{U}, \gamma_i}$. Then a finite \'etale cover $Y \to X$ is tame relative to $K$ if and only if it is tame above each of $\zeta_1,\dots,\zeta_n$ and $\xi_1, \dots, \xi_r$. 
\end{rmk}

\begin{lem}[Good reduction]
    Let $X/K$ be the generic fibre of a connected quasi-compact smooth formal scheme $\fX/K^\circ$.  Assume that $K$ is algebraically closed. 
 Let $\fX_k$ be the special fibre of $\fX$ (not equipped with a log structure!). Then $\pitame(X/K) \simeq \pitame(\fX_k/k)$.
\end{lem}

\begin{proof}
Recall that $\fX_k^{\log}$ denotes the special fibre endowed with log structure. Since smooth implies strictly semistable, we have $\pitame(X/K)\simeq \pitame(\fX_k^{\log}/k)$ by \cref{cor:Abhy-infty-equiv}. Since the log structure of $\fX_k^{\log}$ is constant, by \cref{cor:torus-fibration-tame} and the fact that the stalks of $\overline{\cM}_{\fX^{\log}_k}$ are divisible we get $\pitame(\fX^{\log}_k/k) \simeq \pitame(\fX_k/k)$. 
\end{proof} 

\begin{ex}[Product of polydiscs and circles] \label{ex:polydisc}
    Let $X = \Spa(K\langle x_1, \dots, x_n, y_1^{\pm 1},\dots, y_m^{\pm 1}\rangle)$.   
Then $X$ admits a smooth formal model 
    \[
        \fX = \Spf(K^\circ\langle x_1, \dots, x_n, y_1^{\pm 1},\dots, y_m^{\pm 1}\rangle)
    \]
    with special fibre $Y = (\bA^1)^n\times(\mathbb{G}_m)^m$. Suppose that $K$ is algebraically closed. By the K\"unneth formula for $\pitame(-/k)$ of smooth varieties with snc compactification \cite[Thm. 5.1]{Orgogozo2003:AlterationsGroupeFondamental} and the fact that $\pitame(\bA^1/k)=1$ and $\pitame(\mathbb{G}_m/k)\simeq \widehat{\bZ}'(1)$, we have 
    \[ 
        \pitame(X/K) \simeq \pitame(Y/k)\simeq  \widehat{\bZ}'(1)^m. 
    \]
\end{ex}

\begin{ex}[Annulus]
    Let $X=\Spa(K\langle x, y\rangle/(xy-\pi))$ be the affinoid annulus with outer radius $1$ and inner radius $|\pi|$ for some pseudouniformiser $\pi$.  Assume that $K$ is algebraically closed. 
The formal model $\fX = \Spf(K^\circ\langle x,y\rangle/(xy-\pi))$ is strictly semistable and its special fibre $\fX_k$ is the log scheme considered in \cref{ex:semistable-log-curve,ex:semistable-log-curve-2}. Consequently, we obtain
    \[
        \pitame(X/K)\simeq \pitame(\fX_k/k)\simeq \widehat{\bZ}'(1).
    \]
\end{ex}

\begin{ex}[Punctured disc]
   Let $X = \Spa(K\langle x\rangle)\setminus\{0\}$ be the punctured affinoid disc.  Assume that $K$ is algebraically closed. 
   We can write $X$ as the increasing union of affinoid annuli $X_n = \{|\pi|^n\leq |x|\leq 1\}$. 
   Analytification and restriction yield homomorphisms
   \[
   \pitame(X_n/K) \la \pitame(X_{n+1}/K) \la \ldots \la \pitame(X/K) \la \pitame(\mathbb{G}_{m,K}/K).
   \]
   By the previous example, we have $\pitame(X_n/K) \simeq \widehat{\mathbb{Z}}'(1)$ for every $n$ and the transition maps $\pitame(X_n/K)\to \pitame(X_{n+1}/K)$ as well as the composition of maps to $\pitame(\mathbb{G}_{m,K}/K) \simeq \widehat{\mathbb{Z}}'(1)$ are isomorphisms. In order to show that $\pitame(X/K)$ is also isomorphic to $\widehat{\mathbb{Z}}'(1)$, it remains to show that a connected $Y \in \FEtt(X/K)$ remains connected after restriction to $X_n$ for large enough $n$. As any two points in $Y$ can be joined by a connected quasi-compact subset $Z$ and $Y = \bigcup_n Y_{|X_n}$  is an open covering, this subset $Z$ eventually belongs to some $Y_{|X_n}$ and so these points are in the same connected component of $Y_{|X_n}$ for large $n$. Thus the constant colimit $\varinjlim_n \pi_0(Y_{|X_n})$ injects into $\pi_0(Y) = *$. This concludes the proof. 
\end{ex}

The case $n=3$ of the following computation answers a question of Paul Alexander Helminck.

\begin{ex}[Standard semistable formal scheme] \label{ex:helminck}
    Assume that $K$ is algebraically closed. 
    Let $X = \Spa(K\langle T_1,\ldots,T_n\rangle/(\pi - T_1 \cdot \ldots \cdot T_n))$, for some integer $n \geq 1$ and some pseudouniformiser $\pi \in K$. 
    Then $\fX = \Spf(K^\circ\langle T_1, \ldots, T_n\rangle/(\pi - T_1 \cdot \ldots \cdot T_n))$ is a strictly semistable formal model of~$X$ and by applying 
    \cref{lem:strictly-semistable-formal-scheme-reduces-to-setup364}
    we get the following description of its special fibre 
    \[
        \fX_k = \Spec\left(P_{n}(\pi) \to k[T_1, \ldots,T_n]/(T_1 \cdot \ldots \cdot T_n)\right)
    \]   
    where $P_{n}(\pi) = \Gamma_K^+[e_1, \ldots, e_n]/(e_1 + \ldots + e_n = \nu(\pi))$ as in \cref{lem:strictly-semistable-formal-scheme-reduces-to-setup364}, and $e_i \mapsto T_i$.
    By \cref{cor:Paul} we obtain
    \[
        \pitame(X/K)  = \pitame(\fX_k/k) = \pi_1(P_{n}(\pi))
        = \Hom\big(P_{n}(\pi)^\gp, \widehat{\bZ}'(1)\big) 
        \simeq \widehat{\bZ}'(1)^{n-1}.
    \]    
    Of course \cref{cor:Paul} applies also to more general ``polyhedral affinoids''.
\end{ex}





